\documentclass[11pt,openany,oneside]{book}

\usepackage[T1]{fontenc}
\usepackage{lmodern}
\usepackage{amsmath,amssymb,amsthm,mathtools,mathrsfs}
\usepackage[margin=1.05in]{geometry}
\usepackage{booktabs,tabularx,array,multirow}
\usepackage{enumitem}
\usepackage{graphicx}
\usepackage{tikz}
\usetikzlibrary{arrows.meta,positioning}
\usepackage{pgfplots}
\pgfplotsset{compat=1.18}
\usepackage{microtype}
\usepackage{xcolor}
\usepackage{fancyhdr}
\usepackage[colorlinks=true,linkcolor=black,citecolor=black,urlcolor=blue!55!black]{hyperref}

%========================================================
%  Theorem environments
%========================================================
\theoremstyle{plain}
\newtheorem{theorem}{Theorem}[chapter]
\newtheorem{proposition}[theorem]{Proposition}
\newtheorem{lemma}[theorem]{Lemma}
\newtheorem{corollary}[theorem]{Corollary}

\theoremstyle{definition}
\newtheorem{definition}[theorem]{Definition}
\newtheorem{example}[theorem]{Example}
\newtheorem{exercise}{Exercise}[chapter]

\theoremstyle{remark}
\newtheorem{remark}[theorem]{Remark}

\numberwithin{equation}{chapter}

%========================================================
%  Layout
%========================================================
\linespread{1.08}
\setlength{\parskip}{0.15em}
\setlength{\parindent}{1.4em}

\pagestyle{fancy}
\fancyhf{}
\fancyhead[L]{\small\itshape\nouppercase{\leftmark}}
\fancyhead[R]{\small\thepage}
\renewcommand{\headrulewidth}{0.3pt}
\fancypagestyle{plain}{\fancyhf{}\fancyhead[R]{\small\thepage}\renewcommand{\headrulewidth}{0pt}}

%========================================================
%  Notation
%========================================================
\newcommand{\F}{\mathcal{F}}
\newcommand{\R}{\mathbb{R}}
\newcommand{\C}{\mathbb{C}}
\newcommand{\D}{\mathcal{D}}
\newcommand{\Sw}{\mathcal{S}}
\newcommand{\dd}{\,\mathrm{d}}
\newcommand{\PV}{\operatorname{PV}}
\newcommand{\sgn}{\operatorname{sgn}}
\newcommand{\supp}{\operatorname{supp}}
\newcommand{\ip}[2]{\langle #1,#2\rangle}
\newcommand{\Lone}{L^{1}(\R)}
\newcommand{\Ltwo}{L^{2}(\R)}
\newcommand{\Lonloc}{L^{1}_{\mathrm{loc}}(\R)}
\newcommand{\ep}{\varepsilon}
\newcommand{\fh}{\widehat{f}}
\newcommand{\ph}{\widehat{\varphi}}
\newcommand{\gh}{\widehat{g}}
\newcommand{\hf}{\widehat{H}}

\newcommand{\KeyBox}[1]{%
\begin{center}
\setlength{\fboxsep}{7pt}%
\fbox{$\displaystyle #1$}%
\end{center}}

\setlist[enumerate]{leftmargin=1.6em,itemsep=0.18em,topsep=0.25em}
\setlist[itemize]{leftmargin=1.6em,itemsep=0.18em,topsep=0.25em}

%========================================================
\begin{document}

\begin{titlepage}
\centering
\vspace*{2.4cm}

{\Huge\bfseries The Other Side of the\\[0.3em] Fourier Transform}

\vspace{0.8cm}

{\Large A Study Guide to Distributions,\\[0.15em] Regularization, and Generalized Spectra}

\vspace{0.6cm}

{\large\itshape Every computation is carried out, and every interchange justified.}

\vspace{1.8cm}

\begin{minipage}{0.82\textwidth}
\centering
\rule{\textwidth}{0.4pt}\\[0.9em]
{\large The standing results}\\[1.0em]
$\F[1](\omega)=2\pi\delta(\omega)$\\[0.7em]
$\F[t^{\,n}](\omega)=2\pi i^{\,n}\delta^{(n)}(\omega)$\\[0.7em]
$\F[H](\omega)=\pi\delta(\omega)+\PV\dfrac{1}{i\,\omega}$\\[0.9em]
\rule{\textwidth}{0.4pt}
\end{minipage}

\vfill

\begin{minipage}{0.9\textwidth}
\centering\small
Convention:\quad $\F[f](\omega)=\int_{\R}f(t)e^{-i\omega t}\,\mathrm{d}t$\\[0.35em]
inverse transform normalized by $(2\pi)^{-1}$;\quad no complex conjugation in pairings.
\end{minipage}

\vspace{0.6cm}

{\small Version of 1 November 2026}

\vspace{0.6cm}
\end{titlepage}

\frontmatter

\chapter*{Preface}
\addcontentsline{toc}{chapter}{Preface}

This is a working set of notes, not a textbook and not a script. Its single purpose is to take
the reader from the elementary statement ``the Fourier transform of a constant is a delta'' to a
point where every symbol in that sentence has a definition, every equality has a proof, and every
computation can be redone from scratch.

The material falls into three layers, and the notes keep them visibly separate.

\begin{enumerate}
\item \textbf{Computation.} What the integral is, how it is evaluated, why the answer has the
shape it has. This is the layer one can follow with a first course in calculus.
\item \textbf{Analysis.} Why limits may be interchanged, why a family of kernels has no common
dominating function, why a principal value exists. This layer needs dominated convergence,
Fubini, and integration by parts with boundary terms under control.
\item \textbf{Structure.} What a distribution is, why the space of test functions is chosen as it
is, why the Fourier transform must be moved from $\D'$ to $\Sw'$, and what makes the dual
definition the \emph{only} reasonable one.
\end{enumerate}

The order of the chapters follows the order in which the questions actually arise if one starts
from the constant function and refuses to skip steps: first the failed integral
(Chapter~\ref{ch:constant}), then the objects that appear when one differentiates the answer
(Chapter~\ref{ch:polynomial}), then an ambiguity that forces a decision
(Chapter~\ref{ch:step}), then the analysis behind all the limits used so far
(Chapter~\ref{ch:limits}), then the definition that makes those limits rigorous
(Chapter~\ref{ch:dist}), and finally the extension of the Fourier transform itself
(Chapter~\ref{ch:tespered}).

\subsection*{Prerequisites}

Comfort with $L^{1}$ and $L^{2}$ on the line, uniform convergence, the dominated convergence
theorem, Fubini's theorem, and integration by parts. No functional analysis and no measure theory
beyond the notion of a null set. Two standard facts are used but not proved here, because they
belong to a first course and would derail the narrative: the Riemann--Lebesgue lemma and the
Gaussian integral. Everything else---including the Fourier inversion theorem, in the only form
needed---is proved.

\subsection*{Conventions}

The transform is
\begin{equation}
\F[f](\omega)=\fh(\omega)=\int_{\R}f(t)e^{-i\omega t}\dd t,
\qquad
f(t)=\frac{1}{2\pi}\int_{\R}\fh(\omega)e^{i\omega t}\dd\omega ,
\end{equation}
with $\omega$ the angular frequency and $2\pi$ placed on the inverse transform. All pairings
$\ip{T}{\varphi}$ are complex bilinear: there is no complex conjugation anywhere. A different
normalization changes the constants $2\pi$ in the statements below; it never changes the
mathematics. When comparing with another source, evaluate that source's $\F[1]$ first.

The symbol $\delta$ always denotes the Dirac distribution, $\ip{\delta}{\varphi}=\varphi(0)$;
$H$ is the Heaviside step; $\PV$ is the Cauchy principal value; $\Sw$ and $\Sw'$ are the Schwartz
space and its dual. The word ``distribution'' never means ``probability distribution''.

\subsection*{How to read}

Read a chapter once for the statements, once with a pen to redo the computations, and once to
check the hypotheses of each theorem against the example that motivated it. The exercises in
Appendix~\ref{app:exercises} are chosen so that a correct solution certifies understanding of a
specific step, not of a topic in general.

\begingroup\small\tableofcontents\endgroup

\mainmatter

\chapter{The Fourier Transform of a Constant}
\label{ch:constant}

\begin{quote}
\itshape
A signal that never changes should contain only the frequency $\omega=0$. The claim is
uncontroversial. The integral that is supposed to express it does not converge. This chapter
separates the two facts, computes what can be computed, and identifies exactly what has to be
added to the language in order to write the answer.
\end{quote}

\section{Conventions and the elementary theory}

Throughout,
\begin{equation}
\F[f](\omega)=\fh(\omega)=\int_{\R} f(t)e^{-i\omega t}\dd t ,
\qquad
f(t)=\frac{1}{2\pi}\int_{\R}\fh(\omega)e^{i\omega t}\dd\omega .
\label{eq:ftpair}
\end{equation}
The second formula is not assumed; it is proved in Chapter~\ref{ch:tespered} under hypotheses that
cover everything used here. Until then, any use of it is flagged.

The kernel $e^{-i\omega t}$ is a \emph{frequency probe}. Multiplying a signal by it and integrating
compares the signal with a pure oscillation of frequency $\omega$: when the two agree in frequency
the product is a constant and the contributions add; otherwise they rotate and cancel. This is a
correct description of what the integral \emph{measures}. It is not an argument that the integral
converges, and the distinction is the whole content of the next section.

\begin{proposition}[the $L^{1}$ theory]\label{prop:L1}
Let $f\in\Lone$. Then
\begin{enumerate}
\item the integral in \eqref{eq:ftpair} converges absolutely for every $\omega$, and
\begin{equation}
\|\fh\|_{\infty}\le\|f\|_{1};
\label{eq:L1bound}
\end{equation}
\item $\fh$ is uniformly continuous on $\R$;
\item $\fh(\omega)\to 0$ as $|\omega|\to\infty$ \textup{(Riemann--Lebesgue)}.
\end{enumerate}
\end{proposition}

\begin{proof}
(i) Since $|e^{-i\omega t}|=1$, the integrand is dominated by $|f(t)|$, which does not depend on
$\omega$ and is integrable; \eqref{eq:L1bound} follows by taking suprema. (ii) and (iii) are the
standard consequences of the continuity of translation in $L^{1}$ and of the density of step
functions; both are quoted.
\end{proof}

\begin{example}
$f\equiv 1$ is not in $\Lone$, so Proposition~\ref{prop:L1} says nothing about it. It is worth
insisting on what this does and does not mean: membership in $\Lone$ is a \emph{sufficient}
condition for the transform to exist as an absolutely convergent integral, and its failure is not
a proof that every conceivable interpretation diverges. To decide the question one has to compute.
\end{example}

Two properties of the transform will be used constantly and are recorded once.

\begin{proposition}[elementary symmetries]\label{prop:symmetries}
For any $f$ for which the integrals make sense,
\begin{align}
\F[f(\cdot-a)](\omega)&=e^{-ia\omega}\fh(\omega), \label{eq:timeshift}\\
\F[e^{i\omega_{0}t}f(t)](\omega)&=\fh(\omega-\omega_{0}), \label{eq:modulation}\\
\F[f(a\,\cdot)](\omega)&=\frac{1}{|a|}\,\fh\!\left(\frac{\omega}{a}\right),\qquad a\neq0.
\label{eq:dilation}
\end{align}
If $f$ is real then $\overline{\fh(\omega)}=\fh(-\omega)$; if $f$ is even (odd) then $\fh$ is even
(odd); if $f$ is real and even then $\fh$ is real and even.
\end{proposition}

\begin{proof}
\eqref{eq:timeshift} and \eqref{eq:modulation} are changes of variable in the exponent.
For \eqref{eq:dilation} substitute $t=as$:
\begin{equation}
\int_{\R}f(at)e^{-i\omega t}\dd t
=\frac{1}{|a|}\int_{\R}f(s)e^{-i(\omega/a)s}\dd s
=\frac{1}{|a|}\fh\!\left(\frac{\omega}{a}\right).
\end{equation}
The symmetry statements follow by conjugating and by substituting $t\mapsto-t$.
\end{proof}

\section{The truncated integral, and why it fails}

Let $f\equiv1$. Formally,
\begin{equation}
\fh(\omega)=\int_{\R}e^{-i\omega t}\dd t .
\label{eq:badintegral}
\end{equation}
To give the expression a chance, truncate symmetrically and compute. This is the most favourable
truncation available, so if it fails, all the others fail too.

\begin{lemma}\label{lem:IT}
For $T>0$ and $\omega\in\R$,
\begin{equation}
I_{T}(\omega):=\int_{-T}^{T}e^{-i\omega t}\dd t=
\begin{cases}
\dfrac{2\sin(\omega T)}{\omega}, & \omega\neq0,\\[8pt]
2T, & \omega=0 .
\end{cases}
\label{eq:IT}
\end{equation}
The function $\omega\mapsto I_{T}(\omega)$ is continuous at $\omega=0$.
\end{lemma}

\begin{proof}
For $\omega\neq0$,
\begin{equation}
\int_{-T}^{T}e^{-i\omega t}\dd t
=\left[\frac{e^{-i\omega t}}{-i\omega}\right]_{-T}^{T}
=\frac{e^{-i\omega T}-e^{i\omega T}}{-i\omega}
=\frac{2\sin(\omega T)}{\omega},
\end{equation}
using $\sin x=(e^{ix}-e^{-ix})/(2i)$. For $\omega=0$ the integrand is $1$ and the integral is
$2T$. Continuity at $0$ follows from $\lim_{\omega\to0}\sin(\omega T)/\omega=T$.
\end{proof}

\begin{proposition}\label{prop:nolimit}
Fix $\omega\neq0$. Then $I_{T}(\omega)$ has no limit as $T\to\infty$.
\end{proposition}

\begin{proof}
Put $a=|\omega|>0$. For $n=0,1,2,\dots$ set
\begin{equation}
T_{n}=\frac{\pi/2+2\pi n}{a},
\qquad
S_{n}=\frac{3\pi/2+2\pi n}{a}.
\end{equation}
Then $aT_{n}=\pi/2+2\pi n$ and $aS_{n}=3\pi/2+2\pi n$, so $\sin(aT_{n})=1$ and $\sin(aS_{n})=-1$.
Since $\sin(\omega T)=\pm\sin(aT)$ according to the sign of $\omega$, one gets
\begin{equation}
\frac{2\sin(\omega T_{n})}{\omega}=\pm\frac{2}{a},
\qquad
\frac{2\sin(\omega S_{n})}{\omega}=\mp\frac{2}{a}.
\end{equation}
Both sequences of truncation lengths tend to infinity, while the values stay at two different
numbers. Hence no limit exists.
\end{proof}

At $\omega=0$ the behaviour is different but equally unusable: $I_{T}(0)=2T\to+\infty$. So the
zero frequency does not produce a value either; it produces unbounded growth.

\begin{remark}[three statements, only one of which is correct]
\emph{Wrong:} ``for $\omega\neq0$ the integral equals $0$''. The integral does not exist, so it
cannot equal anything. \emph{Wrong:} ``at $\omega=0$ the transform is infinite'', since $\infty$ is
not a value. \emph{Correct:} for every fixed $\omega\neq0$ the truncated integral oscillates
without limit, and at $\omega=0$ it grows without bound; therefore \eqref{eq:badintegral} is not a
convergent improper integral.
\end{remark}

\begin{remark}[an object worth keeping]\label{rem:sinc}
The function $I_{T}$ is not garbage. It is the Fourier transform of the rectangular window
$w_{T}=\mathbf{1}_{[-T,T]}$, and it will reappear in Chapter~\ref{ch:tespered} as an example of a
family that fails every pointwise test but converges in the distributional sense. The failure
recorded here is a failure of one notion of convergence, not of the object.
\end{remark}

\section{Regularization by a two-sided exponential}

The strategy changes. Instead of evaluating \eqref{eq:badintegral}, we construct a family of
functions that are integrable, computable, and converge to the constant, and we examine where their
transforms go.

For $\ep>0$ set
\begin{equation}
f_{\ep}(t)=e^{-\ep|t|}.
\label{eq:feps}
\end{equation}
Then
\begin{equation}
\|f_{\ep}\|_{1}=\int_{\R}e^{-\ep|t|}\dd t=\frac{2}{\ep}<\infty,
\qquad
f_{\ep}(t)\to1\quad\text{for every fixed }t,
\label{eq:feps-approx}
\end{equation}
so Proposition~\ref{prop:L1} applies to each member of the family. The parameter has a transparent
meaning: $f_{\ep}$ decays on the time scale $1/\ep$, so that as $\ep\downarrow0$ it stays close to
$1$ on ever longer intervals.

\begin{theorem}\label{thm:exptransform}
For every $\ep>0$ and every $\omega\in\R$,
\begin{equation}
\F\big[e^{-\ep|\cdot|}\big](\omega)
=\frac{1}{\ep-i\omega}+\frac{1}{\ep+i\omega}
=\frac{2\ep}{\ep^{2}+\omega^{2}}
=:G_{\ep}(\omega).
\label{eq:Geps}
\end{equation}
\end{theorem}

\begin{proof}
The absolute value changes expression at $t=0$, so split the integral:
\begin{equation}
\int_{\R}e^{-\ep|t|}e^{-i\omega t}\dd t
=\int_{-\infty}^{0}e^{(\ep-i\omega)t}\dd t+\int_{0}^{\infty}e^{-(\ep+i\omega)t}\dd t .
\end{equation}
On the left half line, $|t|=-t$ and the exponent is $(\ep-i\omega)t$, whose real part $\ep$ is
positive; hence
\begin{equation}
\int_{-\infty}^{0}e^{(\ep-i\omega)t}\dd t
=\left[\frac{e^{(\ep-i\omega)t}}{\ep-i\omega}\right]_{-\infty}^{0}
=\frac{1}{\ep-i\omega},
\end{equation}
the lower boundary term vanishing because $|e^{(\ep-i\omega)t}|=e^{\ep t}\to0$ as $t\to-\infty$.
Similarly, on the right half line the real part of $-(\ep+i\omega)$ is $-\ep<0$, so
\begin{equation}
\int_{0}^{\infty}e^{-(\ep+i\omega)t}\dd t
=\left[\frac{e^{-(\ep+i\omega)t}}{-(\ep+i\omega)}\right]_{0}^{\infty}
=\frac{1}{\ep+i\omega}.
\end{equation}
Adding and putting over a common denominator,
\begin{equation}
\frac{1}{\ep-i\omega}+\frac{1}{\ep+i\omega}
=\frac{(\ep+i\omega)+(\ep-i\omega)}{(\ep-i\omega)(\ep+i\omega)}
=\frac{2\ep}{\ep^{2}+\omega^{2}}. \qedhere
\end{equation}
\end{proof}

\begin{remark}[checks that cost nothing]
The computation should be checked against the symmetries of Proposition~\ref{prop:symmetries}.
Since $f_{\ep}$ is real and even, $G_{\ep}$ must be real, even, and nonnegative; it is. The value
at the origin is $G_{\ep}(0)=2/\ep$, which blows up as $\ep\downarrow0$, consistent with the
divergence at zero frequency found in Proposition~\ref{prop:nolimit}. Finally, if $\omega$ is
allowed to be complex, $G_{\ep}$ has poles at $\omega=\pm i\ep$, and these poles approach the real
axis as $\ep\downarrow0$; this is the analytic picture of ``the mass landing on $\omega=0$''. It is
a useful mnemonic, not an argument.
\end{remark}

\section{Peak, width, area}

Three numbers describe the family $G_{\ep}$ completely, and each of them has a definite limit
behaviour.

\begin{proposition}\label{prop:three}
For $\ep>0$,
\begin{equation}
G_{\ep}(0)=\frac{2}{\ep},
\qquad
\operatorname{FWHM}\big(G_{\ep}\big)=2\ep,
\qquad
\int_{\R}G_{\ep}(\omega)\dd\omega=2\pi .
\label{eq:three}
\end{equation}
\end{proposition}

\begin{proof}
The first is immediate. For the second, the half maximum is $\tfrac12 G_{\ep}(0)=1/\ep$, so the
half-maximum points satisfy
\begin{equation}
\frac{2\ep}{\ep^{2}+\omega^{2}}=\frac{1}{\ep}
\iff \ep^{2}+\omega^{2}=2\ep^{2}
\iff |\omega|=\ep ,
\end{equation}
and the full width is $2\ep$. For the third,
\begin{equation}
\int_{\R}\frac{2\ep}{\ep^{2}+\omega^{2}}\dd\omega
=2\ep\cdot\frac{1}{\ep}\left[\arctan\frac{\omega}{\ep}\right]_{-\infty}^{\infty}
=2\left(\frac{\pi}{2}+\frac{\pi}{2}\right)=2\pi . \qedhere
\end{equation}
\end{proof}

Thus the peak grows like $\ep^{-1}$, the width shrinks like $\ep$, and the area is independent of
$\ep$. The product of height and width is constant, which is the quantitative form of the phrase
``spreading in time forces concentration in frequency''. The next proposition makes the mechanism
exact and shows that it is a single identity rather than a collection of coincidences.

\begin{proposition}[dilation]\label{prop:dilation}
Let $g\in\Lone$ and $a>0$. If $f(t)=g(t/a)$ then
\begin{equation}
\fh(\omega)=a\,\widehat{g}(a\omega).
\label{eq:dilate}
\end{equation}
Applied to $g(t)=e^{-|t|}$, whose transform is $2/(1+\omega^{2})$, this gives
\begin{equation}
G_{\ep}(\omega)=\frac{1}{\ep}\cdot\frac{2}{1+(\omega/\ep)^{2}}
=\frac{2\ep}{\ep^{2}+\omega^{2}},
\end{equation}
in agreement with Theorem~\ref{thm:exptransform}.
\end{proposition}

\begin{proof}
Substitute $t=as$, $\dd t=a\dd s$:
\begin{equation}
\fh(\omega)=\int_{\R}g(t/a)e^{-i\omega t}\dd t
=a\int_{\R}g(s)e^{-i\,a\omega\,s}\dd s=a\,\widehat{g}(a\omega). \qedhere
\end{equation}
\end{proof}

\begin{remark}[what the area is not]
The constant value $2\pi$ of the area is a property of this particular family, not a conservation
law. The genuine conservation law is Plancherel's identity
$\|f\|_{2}^{2}=(2\pi)^{-1}\|\fh\|_{2}^{2}$, and it is inapplicable here in the limit: for
$f_{\ep}$,
\begin{equation}
\|f_{\ep}\|_{2}^{2}=\int_{\R}e^{-2\ep|t|}\dd t=\frac{1}{\ep}\longrightarrow\infty ,
\end{equation}
so the constant function is not a finite-energy signal and cannot be treated as an $L^{2}$ limit.
This is worth remembering whenever ``the area of the spectrum'' is mentioned.
\end{remark}

\begin{figure}[t]
\centering
\begin{tikzpicture}
\begin{axis}[width=0.47\textwidth,height=4.3cm,domain=-4:4,samples=200,
  ymin=0,ymax=1.05,xlabel={$t$},ylabel={$f_{\ep}(t)$},
  title={\small time domain},
  legend style={font=\scriptsize,at={(0.03,0.97)},anchor=north west,draw=none},
  ticklabel style={font=\scriptsize},label style={font=\small}]
\addplot[thick,black] {exp(-1*abs(x))};
\addplot[thick,black,dashed] {exp(-0.4*abs(x))};
\legend{$\ep=1$,$\ep=0.4$}
\end{axis}
\end{tikzpicture}\hfill
\begin{tikzpicture}
\begin{axis}[width=0.47\textwidth,height=4.3cm,domain=-4:4,samples=400,
  ymin=0,ymax=5.6,xlabel={$\omega$},ylabel={$G_{\ep}(\omega)$},
  title={\small frequency domain},
  legend style={font=\scriptsize,at={(0.97,0.97)},anchor=north east,draw=none},
  ticklabel style={font=\scriptsize},label style={font=\small}]
\addplot[thick,black] {2*1/(1+x^2)};
\addplot[thick,black,dashed] {2*0.4/(0.16+x^2)};
\legend{$\ep=1$,$\ep=0.4$}
\end{axis}
\end{tikzpicture}
\caption{The same two parameters seen in both domains. The vertical scales differ: the dashed
curve has peak $2/0.4=5$.}
\end{figure}

\section{The normalized kernel and the sampling property}

The family $G_{\ep}$ has total mass $2\pi$. The factor is an artefact of the normalization
\eqref{eq:ftpair} and is removed once and for all.

\begin{definition}\label{def:Keps}
For $\ep>0$ let
\begin{equation}
K_{\ep}(x)=\frac{1}{\pi}\frac{\ep}{x^{2}+\ep^{2}},
\qquad\text{so that}\qquad
K_{\ep}=\frac{G_{\ep}}{2\pi},\qquad \int_{\R}K_{\ep}\dd x=1 .
\label{eq:Keps}
\end{equation}
\end{definition}

\begin{proposition}[mass concentrates at the origin]\label{prop:concentration}
For every $a>0$,
\begin{equation}
\int_{-a}^{a}K_{\ep}(x)\dd x=\frac{2}{\pi}\arctan\frac{a}{\ep}\longrightarrow1
\qquad(\ep\downarrow0).
\label{eq:concentration}
\end{equation}
\end{proposition}

\begin{proof}
Immediate from the antiderivative of $1/(x^{2}+\ep^{2})$; the limit uses
$\arctan y\to\pi/2$ as $y\to\infty$.
\end{proof}

The quantifier order in \eqref{eq:concentration} matters: $a$ is fixed first, and only then does
$\ep$ tend to zero. Equivalently: for every fixed interval around the origin, no matter how short,
the fraction of the total mass inside it tends to one.

\begin{theorem}[sampling]\label{thm:sampling}
Let $\varphi$ be bounded and continuous at $0$. Then
\begin{equation}
\int_{\R}K_{\ep}(x)\varphi(x)\dd x\longrightarrow\varphi(0)
\qquad(\ep\downarrow0).
\label{eq:sampling}
\end{equation}
\end{theorem}

\begin{proof}
Since $\int K_{\ep}=1$, the error is
\begin{equation}
\int_{\R}K_{\ep}(x)\varphi(x)\dd x-\varphi(0)
=\int_{\R}K_{\ep}(x)\big[\varphi(x)-\varphi(0)\big]\dd x .
\end{equation}
Let $\eta>0$. By continuity at $0$ there is $a>0$ with $|\varphi(x)-\varphi(0)|<\eta$ whenever
$|x|<a$. Splitting at $|x|=a$ and using $K_{\ep}\ge0$,
\begin{equation}
\left|\int_{\R}K_{\ep}\big[\varphi(\cdot)-\varphi(0)\big]\right|
\le\eta\int_{|x|<a}K_{\ep}
+2\|\varphi\|_{\infty}\int_{|x|\ge a}K_{\ep}
\le\eta+2\|\varphi\|_{\infty}\Big(1-\frac{2}{\pi}\arctan\frac{a}{\ep}\Big).
\end{equation}
Fix $a$ and let $\ep\downarrow0$: the second term tends to $0$ by
Proposition~\ref{prop:concentration}. Hence
\begin{equation}
\limsup_{\ep\to0}\Big|\int_{\R}K_{\ep}\varphi-\varphi(0)\Big|\le\eta ,
\end{equation}
and $\eta>0$ was arbitrary.
\end{proof}

A second proof, which is often shorter in practice, removes the varying scale by a change of
variable. Substituting $x=\ep u$ in the definition of $K_{\ep}$ gives
\begin{equation}
\int_{\R}K_{\ep}(x)\varphi(x)\dd x
=\frac{1}{\pi}\int_{\R}\frac{\varphi(\ep u)}{1+u^{2}}\dd u ,
\label{eq:substitution-proof}
\end{equation}
and now the weight $(1+u^{2})^{-1}/\pi$ is fixed and the parameter has migrated into the test
function. The integrand is dominated by $\|\varphi\|_{\infty}/(\pi(1+u^{2}))$, which is integrable
and independent of $\ep$; dominated convergence applies directly and gives $\varphi(0)$. The two
proofs are the same estimate seen from two sides, and the second one will be the model for the
general theory in Chapter~\ref{ch:limits}.

\section{The constant, and the delta}

The sampling property \eqref{eq:sampling} is the only thing the kernel family does in the limit.
It is therefore natural to name the limit by what it does.

\begin{definition}[Dirac distribution]\label{def:delta}
The Dirac delta is the rule assigning to each test function $\varphi$ the number $\varphi(0)$. One
writes
\begin{equation}
\ip{\delta}{\varphi}=\varphi(0),
\qquad\text{or}\qquad
\int_{\R}\delta(\omega)\varphi(\omega)\dd\omega=\varphi(0).
\label{eq:deltadef}
\end{equation}
\end{definition}

With this notation, \eqref{eq:sampling} reads $K_{\ep}\to\delta$, and since $G_{\ep}=2\pi K_{\ep}$,
\begin{equation}
G_{\ep}(\omega)=\frac{2\ep}{\ep^{2}+\omega^{2}}\longrightarrow 2\pi\delta(\omega).
\end{equation}
Combining this with \eqref{eq:feps-approx} and with the continuity of the transform in the
appropriate sense (Chapter~\ref{ch:tespered}), one writes
\begin{equation}
\boxed{\;\F[1](\omega)=2\pi\delta(\omega)\;}
\label{eq:F1}
\end{equation}

What \eqref{eq:F1} asserts, precisely, is the following. For every test function $\varphi$,
\begin{equation}
\lim_{\ep\downarrow0}\int_{\R}\F\big[e^{-\ep|\cdot|}\big](\omega)\,\varphi(\omega)\dd\omega
=\lim_{\ep\downarrow0}\int_{\R}G_{\ep}(\omega)\varphi(\omega)\dd\omega
=2\pi\varphi(0)
=\ip{2\pi\delta}{\varphi}.
\end{equation}
The equality is therefore a statement about the action of the transformed family on test
functions. It is not a statement that the divergent integral \eqref{eq:badintegral} has acquired a
value.

\begin{remark}[three consequences worth spelling out]
\begin{enumerate}
\item \textbf{The weight is $2\pi$, not a height.} The number $2\pi$ is the limiting total mass;
the peak height $2/\ep$ is a property of the approximating family and has no limit.
\item \textbf{There is no pointwise interpretation.} $G_{\ep}(\omega)\to0$ for every
$\omega\neq0$, while $G_{\ep}(0)\to\infty$. Neither fact determines the limit; the concentration of
mass does.
\item \textbf{The factor is fixed by the convention.} Had the inverse transform carried
$(2\pi)^{-1/2}$, the constant would have produced $\sqrt{2\pi}\,\delta$. The general rule is
recorded in Appendix~\ref{app:tables}.
\end{enumerate}
\end{remark}

A short independent verification of \eqref{eq:F1} can be given at once and will be simplified in
Chapter~\ref{ch:tespered}. If the transform of $\delta$ is computed first, using the fact that
$\delta$ acts on $e^{-i\omega t}$ by evaluation at $t=0$,
\begin{equation}
\ip{\F\delta}{\varphi}=\ip{\delta}{\ph}=\ph(0)=\int_{\R}\varphi(t)\dd t=\ip{1}{\varphi},
\qquad\text{so}\qquad \F[\delta]=1 .
\end{equation}
Applying the transform twice returns a reflection and a factor $2\pi$
(Theorem~\ref{thm:F2}), which turns $\F[\delta]=1$ into $\F[1]=2\pi\delta$. The two computations
agree, as they must.

\section{Application: the spectrum of a cosine}

The machinery of this chapter answers a question that is usually answered by assertion, namely why a
pure oscillation produces \emph{two} spectral lines.

Let $\omega_{0}\in\R$ and set
\begin{equation}
c_{\ep}(t)=e^{-\ep|t|}\cos(\omega_{0}t).
\end{equation}
Since $c_{\ep}\to\cos(\omega_{0}\cdot)$ pointwise and $|c_{\ep}|\le1$, the family converges in the
same sense as in \eqref{eq:feps-approx}. By
$\cos(\omega_{0}t)=\tfrac12(e^{i\omega_{0}t}+e^{-i\omega_{0}t})$ and the modulation rule
\eqref{eq:modulation},
\begin{equation}
\F[c_{\ep}](\omega)
=\tfrac12\F[f_{\ep}](\omega-\omega_{0})+\tfrac12\F[f_{\ep}](\omega+\omega_{0})
=\tfrac12\Big[G_{\ep}(\omega-\omega_{0})+G_{\ep}(\omega+\omega_{0})\Big].
\label{eq:cos-spectrum}
\end{equation}

\begin{proposition}\label{prop:cosine}
In $\Sw'$,
\begin{equation}
\F[\cos(\omega_{0}\cdot)](\omega)
=\pi\big[\delta(\omega-\omega_{0})+\delta(\omega+\omega_{0})\big].
\label{eq:Fcos}
\end{equation}
\end{proposition}

\begin{proof}
By Theorem~\ref{thm:sampling} applied to $\varphi(\cdot\pm\omega_{0})$,
\begin{equation}
\tfrac12 G_{\ep}(\cdot-\omega_{0})\longrightarrow\pi\delta(\cdot-\omega_{0}),
\qquad
\tfrac12 G_{\ep}(\cdot+\omega_{0})\longrightarrow\pi\delta(\cdot+\omega_{0}),
\end{equation}
and the sum converges to the sum.
\end{proof}

Two checks are instructive. First, putting $\omega_{0}=0$ in \eqref{eq:Fcos} gives
$\pi\delta+\pi\delta=2\pi\delta$, in agreement with \eqref{eq:F1}: the constant is the degenerate case
in which the two lines coincide. Second, the total weight is $2\pi$, distributed as $\pi$ on each of
the two lines; the picture of a single arrow at $+\omega_{0}$ is therefore incomplete, not merely
schematic.

\begin{remark}[why the cosine, and not the sine]
The same computation for $\sin(\omega_{0}t)$ gives
$\tfrac{1}{2i}\big[G_{\ep}(\cdot-\omega_{0})-G_{\ep}(\cdot+\omega_{0})\big]$, hence
$(\pi/i)[\delta(\cdot-\omega_{0})-\delta(\cdot+\omega_{0})]$: two lines of opposite sign. As
$\omega_{0}\to0$ this object tends to $0$, consistently with $\sin(0)=0$. It cannot be used to
approach the constant, which is why the opening question of the chapter is posed with the cosine.
\end{remark}

\section{Summary of the chapter}

\begin{enumerate}
\item For $f\in\Lone$ the transform exists and is bounded by $\|f\|_{1}$
(Proposition~\ref{prop:L1}). The constant is not in $\Lone$, so nothing follows from this.
\item The truncated integral of the constant is $2\sin(\omega T)/\omega$ for $\omega\neq0$ and
$2T$ at $\omega=0$; it has no limit (Lemma~\ref{lem:IT},
Proposition~\ref{prop:nolimit}).
\item Regularization by $e^{-\ep|t|}$ produces the computable spectrum
$2\ep/(\ep^{2}+\omega^{2})$ (Theorem~\ref{thm:exptransform}), with peak $2/\ep$, width $2\ep$, and
constant area $2\pi$ (Proposition~\ref{prop:three}).
\item The normalized kernel $K_{\ep}=G_{\ep}/(2\pi)$ has unit mass, concentrates at the origin, and
samples continuous test functions (Theorem~\ref{thm:sampling}).
\item Writing the limit as $2\pi\delta$ is a definition of the left-hand side of \eqref{eq:F1} as a
functional, not a re-evaluation of a divergent integral.
\end{enumerate}

\chapter{Polynomials and the Derivatives of \texorpdfstring{$\delta$}{delta}}
\label{ch:polynomial}

\begin{quote}
\itshape
If the constant has a spectrum, so does $t$, and so does every polynomial. The answers involve
$\delta'$, $\delta''$, and so on. Before those symbols can mean anything, differentiation has to be
defined for objects that are not functions---and the only definition available is the one that
integration by parts suggests.
\end{quote}

\section{The obstacle, computed}

Take $f(t)=t$. Truncating symmetrically and integrating by parts with $u=t$,
$\dd v=e^{-i\omega t}\dd t$,
\begin{equation}
\int_{-T}^{T}t\,e^{-i\omega t}\dd t
=\left[\frac{t\,e^{-i\omega t}}{-i\omega}\right]_{-T}^{T}
+\frac{1}{i\omega}\int_{-T}^{T}e^{-i\omega t}\dd t
=\frac{2T\cos(\omega T)}{i\omega}+\frac{2\sin(\omega T)}{\omega^{2}} .
\label{eq:t-truncated}
\end{equation}
The first term oscillates with an amplitude that grows linearly in $T$. The failure is therefore
strictly worse than for the constant, where the amplitude was bounded. For $t^{n}$ one gets an
amplitude of order $T^{n}$.

So no amount of truncation will help. The way forward is not analytic but algebraic: obtain the
transform of $t$ from the transform of $1$ by a rule, and postpone the justification of the rule to
Chapter~\ref{ch:tespered}.

\section{The distributional derivative}

\begin{definition}[distributional derivative]\label{def:distderiv}
Let $T$ be a continuous linear functional on $C_{c}^{\infty}(\R)$ (a distribution; see
Chapter~\ref{ch:dist}). Its derivative $T'$ is defined by
\begin{equation}
\ip{T'}{\varphi}:=-\ip{T}{\varphi'},
\qquad
\ip{T^{(n)}}{\varphi}:=(-1)^{n}\ip{T}{\varphi^{(n)}} .
\label{eq:distderiv}
\end{equation}
\end{definition}

The definition is forced by the following observation.

\begin{proposition}[compatibility]\label{prop:compat}
Let $f\in C^{1}(\R)$ with $f'\in\Lonloc$, and let $T_{f}$ denote the distribution
$\varphi\mapsto\int f\varphi$. Then $T_{f}'=T_{f'}$.
\end{proposition}

\begin{proof}
For $\varphi\in C_{c}^{\infty}$, integration by parts on the support of $\varphi$ (which is compact,
so the boundary terms vanish) gives
\begin{equation}
\int_{\R}f'(t)\varphi(t)\dd t=-\int_{\R}f(t)\varphi'(t)\dd t
=-\ip{T_{f}}{\varphi'}=\ip{T_{f}'}{\varphi}. \qedhere
\end{equation}
\end{proof}

Thus \eqref{eq:distderiv} extends the classical derivative wherever the latter exists, and it makes
sense even where no pointwise derivative exists. Two further properties are immediate and will be
used repeatedly: the map $T\mapsto T'$ is linear, and it is continuous with respect to distributional
convergence (Proposition~\ref{prop:derivcont}).

\begin{example}[the derivatives of $\delta$]\label{ex:deltan}
Applying Definition~\ref{def:distderiv} to $\ip{\delta}{\varphi}=\varphi(0)$,
\begin{equation}
\ip{\delta'}{\varphi}=-\ip{\delta}{\varphi'}=-\varphi'(0),
\qquad
\ip{\delta''}{\varphi}=(-1)^{2}\ip{\delta}{\varphi''}=\varphi''(0),
\end{equation}
and in general
\begin{equation}
\ip{\delta^{(n)}}{\varphi}=(-1)^{n}\varphi^{(n)}(0).
\label{eq:deltan}
\end{equation}
So $\delta$ reads off the value at the origin, $\delta'$ reads off (minus) the slope there,
$\delta''$ reads off the curvature. They are different questions about the same point, not
successively sharper peaks.
\end{example}

\begin{remark}[the sign is not decoration]
The factor $(-1)^{n}$ is the boundary term of the integration by parts. Dropping it produces
$\F[t]=-2\pi i\delta'$ instead of $2\pi i\delta'$ and propagates through every formula in this
chapter. It is safer to remember where it comes from than to remember whether it is there.
\end{remark}

\section{Multiplication by \texorpdfstring{$t$}{t} and differentiation in \texorpdfstring{$\omega$}{omega}}

\begin{lemma}[differentiation under the integral]\label{lem:diffint}
Let $f\in\Sw(\R)$. Then $\fh$ is differentiable and
\begin{equation}
\frac{\dd}{\dd\omega}\fh(\omega)=\int_{\R}(-it)f(t)e^{-i\omega t}\dd t .
\label{eq:diffint}
\end{equation}
\end{lemma}

\begin{proof}
Fix $\omega$ and let $h\neq0$. Then
\begin{equation}
\frac{\fh(\omega+h)-\fh(\omega)}{h}
=\int_{\R}f(t)\,\frac{e^{-i(\omega+h)t}-e^{-i\omega t}}{h}\dd t .
\end{equation}
The integrand converges pointwise to $-it\,f(t)e^{-i\omega t}$. To apply dominated convergence,
note that by the mean value theorem the difference quotient is bounded in modulus by
$|t|\,|f(t)|$, which is integrable because $f\in\Sw$. Hence the limit may be taken inside the
integral.
\end{proof}

\begin{theorem}[multiplication by $t$]\label{thm:multt}
For $f\in\Sw$,
\begin{equation}
\F[tf](\omega)=i\,\frac{\dd}{\dd\omega}\fh(\omega).
\label{eq:multt}
\end{equation}
\end{theorem}

\begin{proof}
By Lemma~\ref{lem:diffint},
$\frac{\dd}{\dd\omega}\fh=-i\int_{\R}tf(t)e^{-i\omega t}\dd t=-i\F[tf](\omega)$. Multiplying by
$i$ gives \eqref{eq:multt}.
\end{proof}

\begin{theorem}[differentiation in $t$]\label{thm:derivt}
For $f\in\Sw$,
\begin{equation}
\F[f'](\omega)=i\omega\,\fh(\omega).
\label{eq:derivt}
\end{equation}
\end{theorem}

\begin{proof}
Integrating by parts,
\begin{equation}
\int_{\R}f'(t)e^{-i\omega t}\dd t
=\big[f(t)e^{-i\omega t}\big]_{-\infty}^{\infty}
+i\omega\int_{\R}f(t)e^{-i\omega t}\dd t .
\end{equation}
The boundary term vanishes because $f$ decays faster than any power and $|e^{-i\omega t}|=1$.
\end{proof}

\begin{remark}[the two rules point in opposite directions]
\eqref{eq:multt} multiplies in time and differentiates in frequency; \eqref{eq:derivt}
differentiates in time and multiplies in frequency. Confusing them is the most common error in this
subject. A quick check that catches the confusion: applying the transform to a Gaussian should not
change the Gaussian's decay rate in a way that contradicts either identity.
\end{remark}

\section{From the constant to every polynomial}

\begin{theorem}\label{thm:tn}
For $n=0,1,2,\dots$,
\begin{equation}
\F[t^{\,n}]=2\pi i^{\,n}\delta^{(n)}(\omega).
\label{eq:tn}
\end{equation}
Consequently, for a polynomial $p(t)=\sum_{n=0}^{N}a_{n}t^{n}$,
\begin{equation}
\F[p](\omega)=2\pi\sum_{n=0}^{N}a_{n}i^{\,n}\delta^{(n)}(\omega).
\label{eq:poly}
\end{equation}
\end{theorem}

\begin{proof}
Induction on $n$. For $n=0$ the statement is $\F[1]=2\pi\delta$, which is \eqref{eq:F1}. Suppose
the formula holds for $n$. Then $t^{n+1}=t\cdot t^{n}$, and by \eqref{eq:multt} applied
(distributionally, see \eqref{eq:multt-dist}) to $t^{n}$,
\begin{equation}
\F[t^{n+1}]=i\frac{\dd}{\dd\omega}\F[t^{n}]
=i\frac{\dd}{\dd\omega}\big(2\pi i^{\,n}\delta^{(n)}\big)
=2\pi i^{\,n+1}\delta^{(n+1)} .
\end{equation}
Linearity gives \eqref{eq:poly}.
\end{proof}

The first few cases, written out, are worth having in view:
\begin{equation}
\F[t]=2\pi i\,\delta',\qquad
\F[t^{2}]=-2\pi\,\delta'',\qquad
\F[t^{3}]=-2\pi i\,\delta''',\qquad
\F[t^{4}]=2\pi\,\delta^{(4)} .
\end{equation}
The alternation of the coefficients is exactly the alternation of $i^{\,n}$.

\begin{remark}[what has not been done]
Equation \eqref{eq:tn} is not obtained by substituting a polynomial into a convergent integral;
\eqref{eq:t-truncated} shows that this would fail. It is obtained by an algebraic rule, whose
legitimacy rests on the definition of $\F$ on distributions (Chapter~\ref{ch:tespered}). Until that
chapter, \eqref{eq:tn} is a computation consistent with all the evidence available, and it is flagged
as such.
\end{remark}

\section{Support, and what lives at a single point}

\begin{definition}[support]\label{def:support}
A distribution $T$ \emph{vanishes on an open set} $U$ if $\ip{T}{\varphi}=0$ for every
$\varphi\in C_{c}^{\infty}$ with $\supp\varphi\subset U$. The \emph{support} $\supp T$ is the
complement of the largest open set on which $T$ vanishes.
\end{definition}

\begin{proposition}\label{prop:suppdelta}
$\supp\delta^{(n)}=\{0\}$ for every $n\ge0$.
\end{proposition}

\begin{proof}
If $\varphi$ vanishes in a neighbourhood of $0$, then $\varphi^{(n)}(0)=0$, and by \eqref{eq:deltan}
$\ip{\delta^{(n)}}{\varphi}=0$. Hence $\delta^{(n)}$ vanishes on $\R\setminus\{0\}$, so its support
is contained in $\{0\}$. It is not empty, since $\ip{\delta^{(n)}}{\varphi}\neq0$ for a suitable
$\varphi$ supported near $0$.
\end{proof}

The converse is more informative than the proposition itself: a distribution supported at a single
point must be a finite combination of $\delta$ and its derivatives. This is the precise content of
the phrase ``a polynomial has spectrum concentrated at zero frequency''.

\begin{theorem}[structure of point-supported distributions]\label{thm:support0}
If $T$ is a distribution with $\supp T\subseteq\{0\}$, then there exist an integer $N\ge0$ and
constants $c_{0},\dots,c_{N}$ such that
\begin{equation}
T=\sum_{n=0}^{N}c_{n}\delta^{(n)} .
\end{equation}
\end{theorem}

\begin{proof}
Let $m$ be the order of $T$ near the origin: by the continuity of $T$ there are $C<\infty$ and
$m\in\mathbb{N}$ with
\begin{equation}
|\ip{T}{\psi}|\le C\max_{0\le j\le m}\|\psi^{(j)}\|_{\infty}
\qquad\text{whenever }\supp\psi\subset[-2,2].
\label{eq:order}
\end{equation}
Choose $N\ge m$ and $\chi\in C_{c}^{\infty}$ with $\chi\equiv1$ on $[-1,1]$. For
$\varphi\in C_{c}^{\infty}$, Taylor's formula with integral remainder gives
\begin{equation}
\varphi(x)=\sum_{n=0}^{N}\frac{\varphi^{(n)}(0)}{n!}x^{n}+x^{N+1}\psi(x),
\qquad
\psi(x)=\int_{0}^{1}\frac{(1-s)^{N}}{N!}\varphi^{(N+1)}(sx)\dd s ,
\end{equation}
and $\psi$ is smooth. Multiply by $\chi$ and split:
\begin{equation}
\varphi=\underbrace{\chi\sum_{n=0}^{N}\frac{\varphi^{(n)}(0)}{n!}x^{n}}_{=:P}
+\underbrace{x^{N+1}\psi\chi}_{=:R}.
\end{equation}
Now $P$ is a polynomial vanishing outside $[-2,2]$ after multiplication by $\chi$, and $R$ vanishes
to order $N+1\ge m+1$ at the origin, so by \eqref{eq:order} $\ip{T}{R}=0$. Therefore
\begin{equation}
\ip{T}{\varphi}=\ip{T}{P}
=\sum_{n=0}^{N}\frac{\varphi^{(n)}(0)}{n!}\ip{T}{\chi x^{n}}
=\sum_{n=0}^{N}\tilde c_{n}\,\varphi^{(n)}(0)
=\sum_{n=0}^{N}\big((-1)^{n}\tilde c_{n}\big)\ip{\delta^{(n)}}{\varphi},
\end{equation}
using \eqref{eq:deltan}. Setting $c_{n}=(-1)^{n}\tilde c_{n}$ finishes the proof.
\end{proof}

\begin{corollary}\label{cor:support-poly}
The spectrum of a polynomial is supported at $\{0\}$, and it is a finite combination of
$\delta^{(n)}$; Theorem~\ref{thm:tn} exhibits the coefficients.
\end{corollary}

\begin{remark}[a check by differential equation, and its limits]
If $p$ has degree $N$ then $p^{(N+1)}=0$, so by \eqref{eq:derivt} in its distributional form,
\begin{equation}
(i\omega)^{N+1}\widehat p=0 ,
\qquad\text{i.e.}\qquad
\omega^{N+1}\widehat p=0 .
\end{equation}
Proposition~\ref{prop:keromega} below then shows that $\widehat p$ is a combination of
$\delta,\delta',\dots,\delta^{(N)}$. This recovers the support statement but \emph{not} the
coefficients: the differential equation is compatible with any choice of the $c_{n}$. The
coefficients come from the recursion \eqref{eq:tn}, and only from it.
\end{remark}

\begin{proposition}\label{prop:keromega}
Let $T$ be a distribution with $\omega T=0$, where $\omega T$ denotes multiplication by the
function $\omega\mapsto\omega$. Then $T=c\,\delta$ for some constant $c$.
\end{proposition}

\begin{proof}
Choose $\chi\in C_{c}^{\infty}$ with $\chi\equiv1$ near $0$. For $\varphi\in C_{c}^{\infty}$ define
\begin{equation}
\psi(\omega)=\frac{\varphi(\omega)-\varphi(0)\chi(\omega)}{\omega},
\end{equation}
which is smooth (the numerator vanishes at $0$ together with its first derivative) and compactly
supported. Hence $\varphi=\varphi(0)\chi+\omega\psi$ and
\begin{equation}
\ip{T}{\varphi}=\varphi(0)\ip{T}{\chi}+\ip{\omega T}{\psi}
=c\,\varphi(0)+0=c\ip{\delta}{\varphi},
\end{equation}
with $c=\ip{T}{\chi}$.
\end{proof}

\begin{remark}
The same argument shows why nothing larger can occur: $\omega\delta'=-\delta\neq0$, so a term
$\delta'$ would destroy the relation $\omega T=0$. Only $\delta$ itself is annihilated by
multiplication by $\omega$.
\end{remark}

\section{Smooth approximations to \texorpdfstring{$\delta'$}{delta prime}}

It is useful to see $\delta'$ as a limit of ordinary functions, because the picture explains both
the sign and the failure of ``area'' as a descriptive tool.

Differentiate the Cauchy kernel $K_{\ep}$ of \eqref{eq:Keps}:
\begin{equation}
K_{\ep}'(x)=\frac{1}{\pi}\cdot\frac{-2\ep x}{(x^{2}+\ep^{2})^{2}}
=-\frac{2\ep x}{\pi\,(x^{2}+\ep^{2})^{2}} .
\label{eq:Kprime}
\end{equation}

\begin{proposition}\label{prop:Kprime}
For every $\ep>0$:
\begin{enumerate}
\item $K_{\ep}'$ vanishes only at $0$, is positive for $x<0$ and negative for $x>0$;
\item its extrema occur at $x=\pm\ep/\sqrt3$, where the value is
$\dfrac{9}{8\sqrt3\,\pi\ep^{2}}$;
\item $\displaystyle\int_{\R}K_{\ep}'(x)\dd x=0$;
\item $\displaystyle\int_{\R}|K_{\ep}'(x)|\dd x=\frac{2}{\pi\ep}\to\infty$;
\item for every bounded continuous $\varphi$,
$\ip{K_{\ep}'}{\varphi}=-\ip{K_{\ep}}{\varphi'}\to-\varphi'(0)=\ip{\delta'}{\varphi}$.
\end{enumerate}
\end{proposition}

\begin{proof}
(i) is clear from \eqref{eq:Kprime}. (ii) Differentiating the logarithmic derivative of
$-x/(x^{2}+\ep^{2})^{2}$ gives the condition $\ep^{2}-3x^{2}=0$; substituting yields
$9/(8\sqrt3\pi\ep^{2})$. (iii) $K_{\ep}'$ is the derivative of a function vanishing at $\pm\infty$,
so its integral is $K_{\ep}(\infty)-K_{\ep}(-\infty)=0$. (iv) By symmetry and the substitution
$u=x/\ep$,
\begin{equation}
\int_{\R}|K_{\ep}'|\dd x
=2\int_{0}^{\infty}\frac{2\ep x}{\pi(x^{2}+\ep^{2})^{2}}\dd x
=\frac{4\ep}{\pi}\cdot\frac{1}{2\ep^{2}}
=\frac{2}{\pi\ep}.
\end{equation}
(v) Integration by parts, justified because $K_{\ep}$ is smooth and decays, gives
$\ip{K_{\ep}'}{\varphi}=-\ip{K_{\ep}}{\varphi'}$, and $\varphi'\in C_{b}$ is bounded and continuous,
so Theorem~\ref{thm:sampling} applies to $\varphi'$.
\end{proof}

Part (iii) is the point. A nonnegative kernel representing $\delta$ is characterized by unit mass
concentrating at the origin; there is no analogous description of $\delta'$, because the total
signed mass is zero identically. What survives is the \emph{shape} of the two lobes, and the shape
encodes exactly one number: the derivative of the test function at the origin.

Part (v) also shows that differentiation is compatible with distributional limits: differentiating a
convergent family produces a convergent family. This will be stated generally in
Proposition~\ref{prop:derivcont}.

\section{Worked examples}

\begin{example}\label{ex:poly1}
Compute $\F[3-2t+t^{2}]$. By linearity and Theorem~\ref{thm:tn},
\begin{equation}
\F[3]-2\F[t]+\F[t^{2}]
=3\cdot2\pi\delta-2\cdot 2\pi i\,\delta'-2\pi\,\delta''
=6\pi\delta-4\pi i\,\delta'-2\pi\,\delta'' .
\end{equation}
Consistency check: $3+t^{2}$ is real and even, and contributes the real even part
$6\pi\delta-2\pi\delta''$; $-2t$ is real and odd and contributes the purely imaginary odd part
$-4\pi i\delta'$. This is the symmetry statement of Proposition~\ref{prop:symmetries}, in
distributional form.
\end{example}

\begin{example}\label{ex:delta-prime-test}
Let $\varphi(x)=e^{-x^{2}}+xe^{-x^{2}}$. Then
\begin{equation}
\varphi(0)=1,\qquad
\varphi'(x)=e^{-x^{2}}\big(1-2x-2x^{2}\big),\qquad
\varphi'(0)=1,\qquad
\varphi''(0)=-2 .
\end{equation}
Hence
\begin{equation}
\ip{\delta}{\varphi}=1,\qquad
\ip{\delta'}{\varphi}=-1,\qquad
\ip{\delta''}{\varphi}=-2 .
\end{equation}
\end{example}

\begin{example}\label{ex:ladder}
The identities
\begin{equation}
x\,\delta=0,\qquad
x\,\delta'=-\delta,\qquad
x\,\delta^{(n)}=-n\,\delta^{(n-1)}
\label{eq:xdeltan}
\end{equation}
are proved in Chapter~\ref{ch:dist}. Granting them for a moment, they explain the structure found in
Proposition~\ref{prop:keromega}: multiplication by $\omega$ lowers the order of $\delta^{(n)}$ by
one, so only $\delta^{(0)}$ survives annihilation.
\end{example}

\section{Summary of the chapter}

\begin{enumerate}
\item Polynomials cannot be substituted into \eqref{eq:ftpair}; the truncation grows like $T^{n}$.
\item Differentiation of distributions is defined by transferring the derivative to the test
function, at the price of $(-1)^{n}$ (Definition~\ref{def:distderiv}).
\item Multiplication by $t$ in time corresponds to differentiation in frequency
(Theorem~\ref{thm:multt}); differentiation in time corresponds to multiplication by $i\omega$
(Theorem~\ref{thm:derivt}).
\item The recursion gives $\F[t^{n}]=2\pi i^{n}\delta^{(n)}$
(Theorem~\ref{thm:tn}), and linearity extends it to all polynomials.
\item A distribution supported at the origin is a finite combination of $\delta^{(n)}$
(Theorem~\ref{thm:support0}); this is what ``support at zero frequency'' means.
\item Differentiating the approximating kernel yields $\delta'$ as a distributional limit with two
lobes, zero total mass, and divergent total variation
(Proposition~\ref{prop:Kprime}).
\end{enumerate}

\chapter{The Step Function, the Principal Value, and the Missing \texorpdfstring{$\delta$}{delta}}
\label{ch:step}

\begin{quote}
\itshape
The derivative rules give an equation for the spectrum of the step, and the equation has infinitely
many solutions. The loss of information is not an accident of the method: multiplication by
$\omega$ annihilates $\delta$, exactly as differentiation in time annihilates constants. Recovering
the missing term requires a new object, the Cauchy principal value, and an independent computation.
\end{quote}

\section{The derivative of a step}

Let
\begin{equation}
H(t)=\begin{cases}0,&t<0,\\ 1,& t>0,\end{cases}
\label{eq:H}
\end{equation}
with the value at $0$ left unspecified. The choice at a single point will be irrelevant, and
Chapter~\ref{ch:dist} explains why.

\begin{proposition}\label{prop:Hprime}
As a distribution, $H'=\delta$.
\end{proposition}

\begin{proof}
For $\varphi\in C_{c}^{\infty}$,
\begin{equation}
\ip{H'}{\varphi}=-\ip{H}{\varphi'}=-\int_{0}^{\infty}\varphi'(t)\dd t
=-\big[\varphi(t)\big]_{0}^{\infty}=-\big(0-\varphi(0)\big)=\varphi(0)=\ip{\delta}{\varphi}.
\end{equation}
The boundary at $+\infty$ vanishes because $\varphi$ is compactly supported; the boundary at $0$
produces $\varphi(0)$.
\end{proof}

The mechanism is worth stating separately, since it is the source of every $\delta$ in applications.

\begin{proposition}[piecewise smooth functions]\label{prop:jump}
Let $f$ be $C^{1}$ except at points $a_{j}$ accumulating only at infinity, with jumps
$J_{j}=f(a_{j}^{+})-f(a_{j}^{-})$ and a classical derivative $f'_{\mathrm{cl}}$ off the jump set.
Then
\begin{equation}
Df=f'_{\mathrm{cl}}+\sum_{j}J_{j}\,\delta_{a_{j}},
\label{eq:jump}
\end{equation}
where $\delta_{a}(\varphi)=\varphi(a)$.
\end{proposition}

\begin{proof}
Split the integral $\int f\varphi'$ over the intervals between the jumps and integrate by parts on
each. The interior terms reassemble into $-\int f'_{\mathrm{cl}}\varphi$, and each jump $a_{j}$
contributes $\big(f(a_{j}^{+})-f(a_{j}^{-})\big)\varphi(a_{j})=J_{j}\ip{\delta_{a_{j}}}{\varphi}$.
\end{proof}

\begin{example}\label{ex:abs}
$D|t|=\sgn t$ and $D^{2}|t|=2\delta$. The first has a jump of size $2$ in the derivative of
$|t|$ at the origin; the second differentiates that jump.
\end{example}

\section{Applying the differentiation rule}

By \eqref{eq:derivt}, applied to $H$ in the distributional sense and to $H'=\delta$,
\begin{equation}
i\omega\,\hf(\omega)=\F[H'](\omega)=\F[\delta](\omega)=1 .
\label{eq:iomegaH}
\end{equation}
Read off naively this suggests $\hf(\omega)=1/(i\omega)$. There are two independent reasons why
that cannot be the end of the story.

\begin{enumerate}
\item $1/\omega$ is not locally integrable at $0$: $\int_{0}^{1}\omega^{-1}\dd\omega=\infty$. It
therefore does not define a distribution by integration, and ``divide by $\omega$'' is not an
operation available in this setting.
\item Even granting a meaning to $1/\omega$ near the origin, the equation $i\omega\,\hf=1$ does not
determine $\hf$: if $U$ is one solution then $U+c\delta$ is another, because multiplication by
$\omega$ annihilates $\delta$.
\end{enumerate}

\begin{proposition}\label{prop:omega-delta}
$\omega\delta=0$, that is, $\ip{\omega\delta}{\varphi}=\ip{\delta}{\omega\varphi}=0$ for every
test function $\varphi$.
\end{proposition}

\begin{proof}
By definition of multiplication by a smooth function,
$\ip{\omega\delta}{\varphi}=\ip{\delta}{\omega\varphi}=0\cdot\varphi(0)=0$.
\end{proof}

\begin{remark}[the same loss in both domains]
In the time domain, differentiating kills constants: $(H+C)'=H'$ for every $C$. In the frequency
domain, multiplying by $\omega$ kills $\delta$: $\omega(U+c\delta)=\omega U$ for every $c$. These are
one phenomenon viewed twice. The differentiation rule can therefore only determine the part of the
spectrum away from $\omega=0$; whatever mass sits exactly at zero frequency is invisible to it.
\end{remark}

\section{The principal value}

The function $1/\omega$ has a non-integrable singularity at the origin, but symmetric truncation
removes it. That observation defines a distribution.

\begin{definition}[Cauchy principal value]\label{def:PV}
For $\varphi\in C_{c}^{\infty}(\R)$ set
\begin{equation}
\left\langle\PV\frac{1}{\omega},\varphi\right\rangle
:=\lim_{a\downarrow0}\int_{|\omega|>a}\frac{\varphi(\omega)}{\omega}\dd\omega .
\label{eq:PVdef}
\end{equation}
\end{definition}

The limit can be computed in closed form by pairing the two half lines.

\begin{lemma}\label{lem:PVform}
For every $\varphi\in C_{c}^{\infty}(\R)$ and every $a>0$,
\begin{equation}
\int_{|\omega|>a}\frac{\varphi(\omega)}{\omega}\dd\omega
=\int_{a}^{\infty}\frac{\varphi(x)-\varphi(-x)}{x}\dd x .
\label{eq:PVpaired}
\end{equation}
Consequently
\begin{equation}
\left\langle\PV\frac{1}{\omega},\varphi\right\rangle
=\int_{0}^{\infty}\frac{\varphi(x)-\varphi(-x)}{x}\dd x .
\label{eq:PVint}
\end{equation}
\end{lemma}

\begin{proof}
In the integral over $\omega<-a$ substitute $\omega=-x$ and add the two pieces.
\end{proof}

The point of \eqref{eq:PVpaired} is the numerator: $\varphi(x)-\varphi(-x)$ vanishes at $x=0$ at
first order, so dividing by $x$ leaves a bounded integrand. Each half line separately would
diverge; the pair does not.

\begin{theorem}[existence]\label{thm:PVexists}
The functional \eqref{eq:PVint} is well defined on $C_{c}^{\infty}(\R)$ and is a distribution.
\end{theorem}

\begin{proof}
Let $\varphi\in C_{c}^{\infty}$ and let $K$ be a compact interval containing its support. Split
$\int_{0}^{\infty}$ at $x=1$.

For $0<x\le1$, the mean value theorem gives
\begin{equation}
|\varphi(x)-\varphi(-x)|\le 2x\sup_{|y|\le1}|\varphi'(y)| ,
\qquad\text{hence}\qquad
\left|\frac{\varphi(x)-\varphi(-x)}{x}\right|\le 2\|\varphi'\|_{\infty,[-1,1]} ,
\end{equation}
which is integrable on $(0,1]$.

For $x\ge1$, compact support gives $\varphi(\pm x)=0$ once $x$ exceeds the support radius, so the
integrand vanishes there; in any case $|\varphi(x)-\varphi(-x)|\le 2\|\varphi\|_{\infty}$ and
$1/x$ is integrable on $[1,\infty)$ intersected with the support after bounding.

Linearity is clear. Continuity follows from the displayed estimate together with the bound on
$[1,\infty)$, which together give
\begin{equation}
\left|\left\langle\PV\tfrac{1}{\omega},\varphi\right\rangle\right|
\le C_{K}\Big(\|\varphi'\|_{\infty,[-1,1]}+\|\varphi\|_{\infty}\Big). \qedhere
\end{equation}
\end{proof}

\begin{remark}[the symmetry is part of the definition]\label{rem:PVasym}
It is not true that any way of removing the singularity gives the same answer. Consider the
asymmetric truncation that keeps $a<|\omega|<1$ on the right and $b<|\omega|<1$ on the left, with
$a,b\downarrow0$. By \eqref{eq:PVpaired} and a split at $b$,
\begin{equation}
\int_{a<|\omega|<1}\frac{\varphi(\omega)}{\omega}\dd\omega
=\int_{a}^{1}\frac{\varphi(x)-\varphi(-x)}{x}\dd x+\int_{a}^{b}\frac{\varphi(-x)}{x}\dd x ,
\end{equation}
and the second term equals $\varphi(0)\log(b/a)+o(1)$. Hence the limit exists only when
$b/a\to1$, and then equals the principal value; otherwise the answer differs by
$\varphi(0)\log(b/a)$. The principal value is a specific symmetric limit, not the assertion that
the singularity does not matter.
\end{remark}

Two elementary properties are needed later.

\begin{proposition}\label{prop:PVprops}
$\PV\frac{1}{\omega}$ is an odd distribution, and
\begin{equation}
\omega\,\PV\frac{1}{\omega}=1 .
\label{eq:omegaPV}
\end{equation}
\end{proposition}

\begin{proof}
Oddness follows from \eqref{eq:PVint}: replacing $\varphi$ by its reflection reverses the sign of the
integrand. For \eqref{eq:omegaPV}, use the definition of multiplication by $\omega$ and
\eqref{eq:PVint} with $\omega\varphi$ in place of $\varphi$:
\begin{equation}
\left\langle\omega\PV\tfrac{1}{\omega},\varphi\right\rangle
=\left\langle\PV\tfrac{1}{\omega},\omega\varphi\right\rangle
=\int_{0}^{\infty}\frac{x\varphi(x)-(-x)\varphi(-x)}{x}\dd x
=\int_{0}^{\infty}\big(\varphi(x)+\varphi(-x)\big)\dd x=\int_{\R}\varphi . \qedhere
\end{equation}
\end{proof}

\section{The regularized step, computed}

The natural regularization of $H$ is one-sided: decay only where $H$ is nonzero. Set
\begin{equation}
H_{\ep}(t)=H(t)e^{-\ep t},\qquad \ep>0 .
\label{eq:Heps}
\end{equation}

\begin{theorem}\label{thm:Heps}
For every $\ep>0$ and $\omega\in\R$,
\begin{equation}
\F[H_{\ep}](\omega)=\frac{1}{\ep+i\omega}
=\frac{\ep}{\ep^{2}+\omega^{2}}-i\,\frac{\omega}{\ep^{2}+\omega^{2}} .
\label{eq:Hepssplit}
\end{equation}
\end{theorem}

\begin{proof}
Directly,
\begin{equation}
\int_{0}^{\infty}e^{-(\ep+i\omega)t}\dd t
=\left[\frac{e^{-(\ep+i\omega)t}}{-(\ep+i\omega)}\right]_{0}^{\infty}
=\frac{1}{\ep+i\omega},
\end{equation}
the upper boundary vanishing because $\ep>0$; then rationalize the denominator.
\end{proof}

\begin{remark}[why one-sided]
With the two-sided regularization $e^{-\ep|t|}$ the limit is $1$
(Chapter~\ref{ch:constant}); with the one-sided one it is $H$. The shape of the decay determines the
limit object, not merely the speed of convergence. This is the first place where the choice of
regularization carries information.
\end{remark}

The two terms of \eqref{eq:Hepssplit} have to be treated separately: the real part is even and
concentrates, the imaginary part is odd and must be paired across the origin.

\section{The two limits}

\begin{theorem}[real part]\label{thm:realpart}
\begin{equation}
\frac{\ep}{\ep^{2}+\omega^{2}}\longrightarrow\pi\delta(\omega)
\qquad(\ep\downarrow0),
\label{eq:realpart}
\end{equation}
the limit being taken in the sense of test functions.
\end{theorem}

\begin{proof}
Since $\dfrac{\ep}{\ep^{2}+\omega^{2}}=\pi K_{\ep}$ with $K_{\ep}$ as in \eqref{eq:Keps}, and since
$\ip{K_{\ep}}{\varphi}\to\varphi(0)$ for every bounded continuous $\varphi$
(Theorem~\ref{thm:sampling}), one gets
$\ip{\pi K_{\ep}}{\varphi}\to\pi\varphi(0)=\ip{\pi\delta}{\varphi}$.
\end{proof}

The coefficient $\pi$ is the total mass:
\begin{equation}
\int_{\R}\frac{\ep}{\ep^{2}+\omega^{2}}\dd\omega=\pi ,
\label{eq:masspi}
\end{equation}
independently of $\ep$. Whenever a family of even functions concentrates, its limit is the mass
times $\delta$; this is the single most useful heuristic in the subject, and here it is
exactly what the theorem says.

\begin{theorem}[imaginary part]\label{thm:imagpart}
\begin{equation}
\frac{\omega}{\ep^{2}+\omega^{2}}\longrightarrow\PV\frac{1}{\omega}
\qquad(\ep\downarrow0).
\label{eq:imagpart}
\end{equation}
\end{theorem}

\begin{proof}
Let $\varphi\in C_{c}^{\infty}$. Pair the half lines by substituting $\omega=-x$ in the negative
one:
\begin{equation}
\int_{\R}\frac{\omega}{\ep^{2}+\omega^{2}}\varphi(\omega)\dd\omega
=\int_{0}^{\infty}\frac{x}{\ep^{2}+x^{2}}\big[\varphi(x)-\varphi(-x)\big]\dd x .
\label{eq:imagsplit}
\end{equation}
We apply dominated convergence to \eqref{eq:imagsplit} and must exhibit a dominating function
independent of $\ep$.

\emph{Near the origin}, for $0<x\le1$, use $x/(\ep^{2}+x^{2})\le1/x$ together with the mean value
estimate $|\varphi(x)-\varphi(-x)|\le2x\sup_{|y|\le1}|\varphi'|$:
\begin{equation}
\left|\frac{x}{\ep^{2}+x^{2}}\big[\varphi(x)-\varphi(-x)\big]\right|
\le 2\sup_{|y|\le1}|\varphi'(y)| ,
\end{equation}
a constant, hence integrable on $(0,1]$.

\emph{Far from the origin}, for $x\ge1$, the same inequality $x/(\ep^{2}+x^{2})\le1/x$ combined
with compact support gives a bound $C/x$ on a bounded interval, hence integrable on $[1,\infty)$.

For each fixed $x>0$, $x/(\ep^{2}+x^{2})\to1/x$ as $\ep\downarrow0$. Dominated convergence in
\eqref{eq:imagsplit} therefore yields
\begin{equation}
\int_{\R}\frac{\omega}{\ep^{2}+\omega^{2}}\varphi(\omega)\dd\omega
\longrightarrow\int_{0}^{\infty}\frac{\varphi(x)-\varphi(-x)}{x}\dd x
=\left\langle\PV\frac{1}{\omega},\varphi\right\rangle ,
\end{equation}
which is \eqref{eq:imagpart}.
\end{proof}

\begin{remark}[why the two regions must be separated]
The kernel $\omega/(\ep^{2}+\omega^{2})$ is small near the origin and small at infinity, with a
maximum of $1/(2\ep)$ at $\omega=\pm\ep$. A single global estimate cannot capture both ends; the
split at $x=1$ makes the near part bounded by the oscillation of $\varphi$ and the far part bounded
by the decay of $\varphi$. This ``control near the singular point by smoothness, control at
infinity by decay'' is the standard pattern, and it recurs verbatim in
Chapter~\ref{ch:limits}.
\end{remark}

\section{The spectrum of the step}

Combining Theorems~\ref{thm:Heps}, \ref{thm:realpart} and \ref{thm:imagpart}, and using
$1/i=-i$,
\begin{equation}
\F[H](\omega)=\pi\delta(\omega)-i\,\PV\frac{1}{\omega}
=\pi\delta(\omega)+\PV\frac{1}{i\omega} .
\label{eq:FH}
\end{equation}

The coefficient $\pi$ in front of $\delta$ deserves an independent verification, since it was
obtained from a computation rather than from a structural argument.

\begin{theorem}[symmetry determines the coefficient]\label{thm:symcoef}
Let $c\in\C$ and suppose $\F[H]=c\,\delta-i\,\PV\frac{1}{\omega}$. Then $c=\pi$.
\end{theorem}

\begin{proof}
As distributions, $H(t)+H(-t)=1$: the identity holds for every $t\neq0$, and the value at a single
point does not affect the distribution. Applying $\F$ and using
$\F[f(-\cdot)](\omega)=\fh(-\omega)$,
\begin{equation}
\hf(\omega)+\hf(-\omega)=2\pi\delta(\omega).
\label{eq:Hsym}
\end{equation}
Substituting the ansatz and using that $\PV\frac{1}{\omega}$ is odd while $\delta$ is even,
\begin{equation}
\big(c\delta-i\PV\tfrac{1}{\omega}\big)+\big(c\delta+i\PV\tfrac{1}{\omega}\big)=2c\,\delta .
\end{equation}
Comparing with \eqref{eq:Hsym} gives $2c=2\pi$, hence $c=\pi$.
\end{proof}

\begin{corollary}[the sign function]\label{cor:sgn}
\begin{equation}
\F[\sgn t](\omega)=-2i\,\PV\frac{1}{\omega}=\frac{2}{i}\,\PV\frac{1}{\omega} .
\label{eq:sgn}
\end{equation}
\end{corollary}

\begin{proof}
Since $H=\tfrac12+\tfrac12\sgn$ and $\F[1]=2\pi\delta$,
\begin{equation}
\F[\sgn t]=2\Big(\F[H]-\tfrac12\F[1]\Big)
=2\big(\pi\delta-i\PV\tfrac{1}{\omega}-\pi\delta\big)
=-2i\,\PV\tfrac{1}{\omega}. \qedhere
\end{equation}
\end{proof}

\begin{remark}[why no $\delta'$ can be added]
The equation $i\omega\hf=1$ is compatible with adding $c\delta$, since $\omega\delta=0$, but not
with adding $\delta'$, since $\omega\delta'=-\delta\neq0$
(\eqref{eq:xdeltan}). Proposition~\ref{prop:keromega} makes this precise: the solutions of
$\omega T=1$ differ exactly by multiples of $\delta$. So \eqref{eq:FH} is the general solution, and
the only question was the value of $c$, settled by \eqref{eq:Hsym}.
\end{remark}

\section{Boundary values}

The computation \eqref{eq:Hepssplit} also produces the standard notation of mathematical physics.
For $\ep\downarrow0$,
\begin{equation}
\frac{1}{x\pm i0}:=\lim_{\ep\downarrow0}\frac{1}{x\pm i\ep}
=\PV\frac{1}{x}\mp i\pi\,\delta(x),
\label{eq:boundary}
\end{equation}
where the limit is distributional. The notation is a shorthand for the limit, not a substitution of
a complex zero.

To see the sign in \eqref{eq:FH} by this route, note that
\begin{equation}
\frac{1}{\ep+i\omega}=\frac{-i}{\omega-i\ep}\longrightarrow\frac{-i}{\omega-i0}
=-i\Big(\PV\frac{1}{\omega}+i\pi\delta\Big)=\pi\delta-i\,\PV\frac{1}{\omega},
\end{equation}
which reproduces \eqref{eq:FH}. The two minus signs---one from $1/i=-i$, one from the lower sign in
\eqref{eq:boundary}---cancel in the $\delta$ term and add in the $\PV$ term; this is the most
common place for a sign error to enter, and it is worth doing the computation twice.

\section{The transform of the principal value}

One more table entry can be obtained for free from what has already been proved, and it illustrates how
the reflection identity \eqref{eq:F2} is used in practice.

\begin{proposition}\label{prop:PVtransform}
In $\Sw'$,
\begin{equation}
\F\Big[\PV\frac{1}{t}\Big](\omega)=-i\pi\,\sgn\omega .
\label{eq:PVtransform}
\end{equation}
\end{proposition}

\begin{proof}
Apply $\F^{2}T=2\pi\,T(-\cdot)$ to $T=\sgn t$. By \eqref{eq:sgn},
$\F[\sgn t]=-2i\PV(1/\omega)$, so
\begin{equation}
\F\Big[\F[\sgn t]\Big]=\F\Big[-2i\,\PV\frac{1}{\omega}\Big]=-2i\,\F\Big[\PV\frac{1}{\omega}\Big]
\end{equation}
while the reflection identity gives
$\F^{2}[\sgn t]=2\pi\,\sgn(-t)=-2\pi\,\sgn t$. Equating,
\begin{equation}
-2i\,\F\Big[\PV\frac{1}{\omega}\Big]=-2\pi\,\sgn\omega ,
\qquad\text{hence}\qquad
\F\Big[\PV\frac{1}{\omega}\Big]=\frac{\pi}{i}\,\sgn\omega=-i\pi\,\sgn\omega . \qedhere
\end{equation}
\end{proof}

\begin{remark}[a consistency check on the factor]
The answer must be odd and purely imaginary, because $\PV(1/t)$ is odd and real; it is. A sharper
check is to apply $\F$ once more and use \eqref{eq:F2}. On one hand
\begin{equation}
\F\Big[\F\Big[\PV\frac{1}{t}\Big]\Big]=\F\big[-i\pi\,\sgn\omega\big]=-i\pi\cdot\big(-2i\big)\PV\frac{1}{\omega}
=-2\pi\,\PV\frac{1}{\omega},
\end{equation}
using \eqref{eq:sgn}; on the other hand \eqref{eq:F2} gives
$2\pi\,\PV\frac{1}{-t}=-2\pi\,\PV\frac{1}{t}$. The two agree, as they must, and the factor
$\pi$ is consistent with \eqref{eq:Fcos}.
\end{remark}

\section{Summary of the chapter}

\begin{enumerate}
\item $H'=\delta$, and more generally each jump of size $J$ contributes $J\delta$ to the
distributional derivative (Propositions~\ref{prop:Hprime}, \ref{prop:jump}).
\item The differentiation rule yields $i\omega\hf=1$, which does not determine $\hf$ because
$\omega\delta=0$ (Proposition~\ref{prop:omega-delta}).
\item $1/\omega$ is not locally integrable; its symmetric regularization is the principal value,
which is a well-defined distribution (Theorems~\ref{thm:PVexists}) with
$\omega\PV\frac{1}{\omega}=1$.
\item The asymmetric truncation shows that the symmetry in the definition is essential
(Remark~\ref{rem:PVasym}).
\item Regularizing the step one-sidedly and separating real and imaginary parts gives
$\pi\delta$ and $-i\PV\frac{1}{\omega}$ (Theorems~\ref{thm:realpart}, \ref{thm:imagpart}).
\item The coefficient $\pi$ is forced by $H(t)+H(-t)=1$ (Theorem~\ref{thm:symcoef}).
\item Consequently $\F[H]=\pi\delta+\PV\frac{1}{i\omega}$ and
$\F[\sgn t]=-2i\PV\frac{1}{\omega}$.
\end{enumerate}

\chapter{Limits, Integrals, and Approximate Identities}
\label{ch:limits}

\begin{quote}
\itshape
Everything so far has used the word ``tends'' without saying in which sense. This chapter supplies
the analysis: three inequivalent notions of convergence, one theorem that licenses the interchange
of limit and integral, and a complete proof that any kernel family with unit mass concentrating at a
point samples continuous functions at that point.
\end{quote}

\section{Three modes of convergence}

Let $f_{n},f$ be functions on $\R$ and let $T_{n},T$ be distributions.

\begin{definition}\label{def:modes}
\begin{enumerate}
\item $f_{n}\to f$ \emph{pointwise} if $f_{n}(x)\to f(x)$ for each fixed $x$.
\item $f_{n}\to f$ \emph{uniformly} if $\sup_{x}|f_{n}(x)-f(x)|\to0$.
\item $T_{n}\to T$ \emph{in the distributional sense} if $\ip{T_{n}}{\varphi}\to\ip{T}{\varphi}$ for
each fixed test function $\varphi$.
\end{enumerate}
\end{definition}

The three are not ordered by strength in the way one might hope, and on an unbounded domain none of
them controls the behaviour of integrals.

\begin{example}[uniform convergence does not preserve the integral]\label{ex:uniform}
Let $f_{n}(x)=n^{-1}\mathbf{1}_{[0,n]}(x)$. Then $\sup_{x}|f_{n}(x)|=1/n\to0$, so $f_{n}\to0$
uniformly, while
\begin{equation}
\int_{\R}f_{n}(x)\dd x=\frac{1}{n}\cdot n=1
\qquad\text{for every }n .
\end{equation}
The functions become shorter and longer at the same rate; the area is constant.
\end{example}

\begin{example}[pointwise convergence is compatible with a delta limit]\label{ex:delta-limit}
Let $g_{n}(x)=n\,\mathbf{1}_{(0,1/n)}(x)$. For fixed $x\neq0$ one has $g_{n}(x)=0$ as soon as
$n>1/x$, so $g_{n}\to0$ pointwise. Nevertheless, for $\varphi$ continuous at $0$,
\begin{equation}
\ip{g_{n}}{\varphi}=n\int_{0}^{1/n}\varphi(x)\dd x\longrightarrow\varphi(0),
\end{equation}
that is, $g_{n}\to\delta$ distributionally. The same sequence tends to $0$ pointwise and to
$\delta$ in the distributional sense; the two statements are about different objects.
\end{example}

\begin{remark}
For the Cauchy kernel the pointwise statement has to be qualified once more: $K_{\ep}(x)\to0$ for
$x\neq0$, but $K_{\ep}(0)=1/(\pi\ep)\to\infty$. Statements of the form ``the kernel tends to zero
everywhere'' are therefore false, and the exception at the origin is precisely the carrier of the
mass.
\end{remark}

\section{When may a limit enter an integral?}

\begin{theorem}[dominated convergence]\label{thm:DCT}
Let $f_{n}\to f$ almost everywhere. Suppose there is a single $g\in\Lone$ with
$|f_{n}(x)|\le g(x)$ almost everywhere, for all $n$. Then
\begin{equation}
\int_{\R}f_{n}\longrightarrow\int_{\R}f .
\end{equation}
\end{theorem}

The hypothesis that is usually misread is the phrase \emph{a single} $g$: a different bound for each
$n$ is not enough.

\begin{proposition}[no common dominator for the Cauchy kernels]\label{prop:nodominator}
There is no $g\in\Lone$ with $K_{\ep}\le g$ for all $\ep\in(0,1)$.
\end{proposition}

\begin{proof}
Suppose such a $g$ exists. Since $K_{\ep}(x)\to0$ for each $x\neq0$, dominated convergence gives
$\int K_{\ep}\to0$. But $\int K_{\ep}=1$ for every $\ep$, a contradiction.
\end{proof}

\begin{remark}[what the test function buys]
The family $e^{-\ep|t|}$ is dominated by the constant $1$, which is not integrable; so
Theorem~\ref{thm:DCT} does not apply to $\int e^{-\ep|t|}\dd t$ over the whole line. Multiply by a
test function, however, and
\begin{equation}
\big|e^{-\ep|t|}\varphi(t)\big|\le|\varphi(t)|\in\Lone ,
\end{equation}
a bound independent of $\ep$ and integrable. Dominated convergence then gives
$\ip{e^{-\ep|\cdot|}}{\varphi}\to\ip{1}{\varphi}$ for every $\varphi\in\Sw$.

The introduction of test functions is therefore not a formality: it is the operation that replaces a
non-integrable dominating function by an integrable one. This single observation is the technical
reason the whole theory is built on pairings.
\end{remark}

\section{Approximate identities}

\begin{definition}\label{def:approxid}
Let $k\in\Lone$ with $\int_{\R}k=1$. For $\ep>0$ put
\begin{equation}
k_{\ep}(x)=\frac{1}{\ep}\,k\!\left(\frac{x}{\ep}\right).
\label{eq:dilate-k}
\end{equation}
The family $(k_{\ep})_{\ep>0}$ is called an \emph{approximate identity}.
\end{definition}

\begin{lemma}[mass and tails]\label{lem:mass-tail}
For every $\ep>0$, $\int_{\R}k_{\ep}=1$. For every $a>0$,
\begin{equation}
\int_{|x|>a}k_{\ep}(x)\dd x=\int_{|u|>a/\ep}k(u)\dd u\longrightarrow0
\qquad(\ep\downarrow0).
\label{eq:tail}
\end{equation}
\end{lemma}

\begin{proof}
The first statement is the substitution $x=\ep u$. For the second, the same substitution converts the
left side into the tail of the fixed integrable function $k$ over the expanding set
$\{|u|>a/\ep\}$; the tails of an integrable function vanish.
\end{proof}

\begin{theorem}[sampling]\label{thm:sampling-general}
Let $(k_{\ep})$ be an approximate identity and let $\varphi$ be bounded and continuous at $0$. Then
\begin{equation}
\int_{\R}k_{\ep}(x)\varphi(x)\dd x\longrightarrow\varphi(0).
\label{eq:sampling-general}
\end{equation}
\end{theorem}

\begin{proof}
Because the total mass is $1$, the difference is
\begin{equation}
\int_{\R}k_{\ep}\varphi-\varphi(0)=\int_{\R}k_{\ep}(x)\big[\varphi(x)-\varphi(0)\big]\dd x .
\end{equation}
Let $\eta>0$. Continuity at $0$ gives $a>0$ such that $|x|<a$ implies
$|\varphi(x)-\varphi(0)|<\eta$. Splitting at $|x|=a$ and using $k_{\ep}\ge0$ and
$\|k_{\ep}\|_{1}=1$,
\begin{equation}
\left|\int_{\R}k_{\ep}\varphi-\varphi(0)\right|
\le\eta\int_{|x|<a}k_{\ep}+\int_{|x|\ge a}k_{\ep}\,2\|\varphi\|_{\infty}
\le\eta+2\|\varphi\|_{\infty}\int_{|x|\ge a}k_{\ep}.
\end{equation}
Now fix $a$ and let $\ep\downarrow0$; the last term tends to $0$ by \eqref{eq:tail}. Hence the
superior limit is at most $\eta$, for every $\eta>0$.
\end{proof}

\begin{remark}[the structure of the estimate]
The proof uses two different mechanisms in two different regions: near the origin the smallness of
$\varphi(x)-\varphi(0)$ comes from \emph{continuity}, and far from the origin the smallness of the
remaining mass comes from \emph{integrability}. The parameter $a$ is fixed before $\ep$ is sent to
zero, and $\eta$ is sent to zero last. Reversing this order of quantifiers is the standard way to
produce a wrong proof.
\end{remark}

For the specific Cauchy kernel a change of variable gives a shorter argument, and the comparison is
instructive.

\begin{proposition}\label{prop:cauchy-sub}
For bounded continuous $\varphi$,
\begin{equation}
\int_{\R}K_{\ep}(x)\varphi(x)\dd x=\frac{1}{\pi}\int_{\R}\frac{\varphi(\ep u)}{1+u^{2}}\dd u
\longrightarrow\varphi(0).
\end{equation}
\end{proposition}

\begin{proof}
Substitute $x=\ep u$. The integrand is dominated by
$\|\varphi\|_{\infty}/(\pi(1+u^{2}))$, which does not involve $\ep$ and is integrable; for each fixed
$u$, $\varphi(\ep u)\to\varphi(0)$. Dominated convergence applies, and
$\frac{1}{\pi}\int(1+u^{2})^{-1}\dd u=1$.
\end{proof}

The two proofs differ in what is held fixed. In Theorem~\ref{thm:sampling-general} the kernel varies
and the test function is fixed; in Proposition~\ref{prop:cauchy-sub} the kernel is rescaled into a
fixed weight and the variation is transferred to the test function. Only the second one produces a
dominating function, which is why the first one needs the two-region estimate instead.

\section{Convolution and the identity element}

\begin{definition}
For suitable $k,f$,
\begin{equation}
(k*f)(x)=\int_{\R}k(x-y)f(y)\dd y=\int_{\R}k(y)f(x-y)\dd y .
\end{equation}
\end{definition}

\begin{theorem}\label{thm:conv}
Let $(k_{\ep})$ be an approximate identity and let $f$ be bounded and continuous at $x$. Then
$(k_{\ep}*f)(x)\to f(x)$. If $f$ is bounded and uniformly continuous, the convergence is uniform;
if $f\in L^{p}(\R)$ with $1\le p<\infty$, then $k_{\ep}*f\to f$ in $L^{p}$.
\end{theorem}

\begin{proof}[Proof of the pointwise statement]
Apply Theorem~\ref{thm:sampling-general} to the function $y\mapsto f(x-y)$, which is bounded and
continuous at $y=0$:
\begin{equation}
(k_{\ep}*f)(x)-f(x)=\int_{\R}k_{\ep}(y)\big[f(x-y)-f(x)\big]\dd y\longrightarrow0 . \qedhere
\end{equation}
\end{proof}

\begin{proposition}[$\delta$ is the identity]\label{prop:delta-identity}
Formally and, once $\delta$ is a distribution, rigorously,
\begin{equation}
(\delta*f)(x)=\int_{\R}\delta(y)f(x-y)\dd y=f(x).
\end{equation}
\end{proposition}

\begin{proof}
By the defining property of $\delta$ applied to the test function $y\mapsto f(x-y)$.
\end{proof}

So the phrase ``approximate identity'' is literal: convolution with $k_{\ep}$ approximates the
identity map on functions, and $\delta$ is that map.

\section{How much of the structure is needed}

\begin{proposition}[general sufficient conditions]\label{prop:general-kernels}
Let $(k_{\ep})_{\ep>0}\subset\Lone$ satisfy
\begin{enumerate}
\item $\displaystyle\int_{\R}k_{\ep}=1$ for every $\ep$;
\item $\displaystyle\sup_{\ep}\|k_{\ep}\|_{1}<\infty$;
\item $\displaystyle\int_{|x|>a}|k_{\ep}(x)|\dd x\to0$ for every $a>0$.
\end{enumerate}
Then $\ip{k_{\ep}}{\varphi}\to\varphi(0)$ for every bounded continuous $\varphi$.
\end{proposition}

\begin{proof}
Repeat the proof of Theorem~\ref{thm:sampling-general} with $k_{\ep}$ replaced by $|k_{\ep}|$; the
near term is bounded by $\eta\sup_{\ep}\|k_{\ep}\|_{1}$ and the far term by
$2\|\varphi\|_{\infty}\int_{|x|\ge a}|k_{\ep}|$, which tends to $0$.
\end{proof}

Nonnegativity is thus a convenience, not a necessity.

\begin{proposition}[mass can escape]\label{prop:escaping}
Let $k_{n}=\tfrac12\mathbf{1}_{[n-1,n+1]}$. Then $\int_{\R}k_{n}=1$ for every $n$, but
$\ip{k_{n}}{\varphi}\to0$ for every $\varphi\in\Sw$; that is, $k_{n}\to0$, not $\delta$.
\end{proposition}

\begin{proof}
For $\varphi\in\Sw$ one has $\sup_{|x-n|\le1}|\varphi(x)|=O(n^{-2})$ by rapid decay, hence
\begin{equation}
|\ip{k_{n}}{\varphi}|\le\frac{1}{2}\int_{n-1}^{n+1}|\varphi(x)|\dd x=O(n^{-2})\longrightarrow0 .
\end{equation}
\end{proof}

\begin{remark}[why this does not contradict mass conservation]
The total mass is conserved, but it escapes to infinity. In the topology of $\Sw'$ the only
continuous tests are against rapidly decaying functions, and the constant function $1$ is not one of
them. Consequently ``total mass'' is not a continuous functional on $\Sw'$, and a family may conserve
mass while converging to zero. This is not a pathology to be avoided; it is the reason the theory
must be formulated with test functions rather than with total integrals.
\end{remark}

\begin{remark}[the sinc kernels are outside this theorem]
The kernels
\begin{equation}
\widehat{w_{T}}(\omega)=\frac{2\sin(T\omega)}{\omega},
\qquad w_{T}=\mathbf{1}_{[-T,T]},
\end{equation}
have total mass $2\pi$ and converge to $2\pi\delta$, but they satisfy none of the hypotheses above:
they change sign infinitely often, and $\int|\widehat{w_{T}}|$ grows like $\log T$. Their convergence
is proved in Chapter~\ref{ch:tespered} by a different route, using the continuity of $\F$ on
$\Sw'$. The existence of this example is the reason Proposition~\ref{prop:general-kernels} is stated
as a sufficient, not a necessary, condition.
\end{remark}

\section{Summary of the chapter}

\begin{center}
\small
\begin{tabular}{@{}p{3.2cm}p{9.6cm}@{}}
\toprule
Statement & Content \\
\midrule
Three modes & pointwise, uniform, distributional; none implies control of integrals
(Examples~\ref{ex:uniform}, \ref{ex:delta-limit}) \\
Licence & dominated convergence, with a \emph{single} integrable bound
(Theorem~\ref{thm:DCT}) \\
Obstruction & no integrable bound exists for $K_{\ep}$
(Proposition~\ref{prop:nodominator}) \\
Sampling & unit mass plus tail concentration implies $\ip{k_{\ep}}{\varphi}\to\varphi(0)$
(Theorem~\ref{thm:sampling-general}) \\
Identity & $\delta*f=f$ (Proposition~\ref{prop:delta-identity}) \\
Necessity & nonnegativity is optional, concentration is not
(Proposition~\ref{prop:general-kernels}) \\
Escape & unit mass alone allows convergence to $0$
(Proposition~\ref{prop:escaping}) \\
\bottomrule
\end{tabular}
\end{center}

\chapter{Distributions}
\label{ch:dist}

\begin{quote}
\itshape
An object may be described by what it does to probes rather than by what value it takes at each
point. Once that shift is made, $\delta$ stops being a paradox: it is a perfectly ordinary element of
a space whose definition requires nothing more than smoothness and compact support of the probes.
\end{quote}

\section{Test functions}

\begin{definition}\label{def:D}
The space of test functions is
\begin{equation}
\D(\R)=C_{c}^{\infty}(\R)=\{\varphi:\R\to\C \;:\; \varphi\in C^{\infty},\ \supp\varphi
\text{ compact}\}.
\end{equation}
\end{definition}

The two requirements do different work. Smoothness allows derivatives of arbitrary order to be
transferred from the object to the probe, which is how differentiation will be defined. Compact
support makes boundary terms vanish in every integration by parts, and it confines the observation to
a bounded region, which is what allows $\Lonloc$ functions to act.

\begin{example}[a smooth function with compact support]\label{ex:bump}
\begin{equation}
\eta(x)=\begin{cases}\exp\!\left(-\dfrac{1}{1-x^{2}}\right),&|x|<1,\\[8pt] 0,&|x|\ge1 .\end{cases}
\end{equation}
As $x\to\pm1$ the exponent tends to $-\infty$ faster than any power, so $\eta$ and all its
derivatives tend to $0$ at the boundary, and the piecewise definition is $C^{\infty}$. The function
is positive on $(-1,1)$ and vanishes identically outside.
\end{example}

\begin{example}
$e^{-x^{2}}$ is smooth and decays rapidly, but its support is all of $\R$; it belongs to the Schwartz
space of Chapter~\ref{ch:tespered} and not to $\D$.
\end{example}

\begin{definition}[convergence in $\D$]
$\varphi_{j}\to\varphi$ in $\D$ means: there is a compact $K$ containing every $\supp\varphi_{j}$ and
$\supp\varphi$, and for every $k\ge0$,
\begin{equation}
\sup_{x\in K}\big|\varphi_{j}^{(k)}(x)-\varphi^{(k)}(x)\big|\longrightarrow0 .
\end{equation}
\end{definition}

The common compact set is part of the definition. A sequence of probes that shrinks while moving off
to infinity is not convergent in $\D$, and admitting such sequences would destroy the continuity
properties the theory needs.

\section{Distributions}

\begin{definition}\label{def:distribution}
A \emph{distribution} on $\R$ is a map $T:\D(\R)\to\C$, written $\varphi\mapsto\ip{T}{\varphi}$,
which is
\begin{enumerate}
\item \emph{linear}: $\ip{T}{a\varphi+b\psi}=a\ip{T}{\varphi}+b\ip{T}{\psi}$ for $a,b\in\C$;
\item \emph{continuous}: $\varphi_{j}\to\varphi$ in $\D$ implies $\ip{T}{\varphi_{j}}\to\ip{T}{\varphi}$.
\end{enumerate}
The space of distributions is denoted $\D'(\R)$.
\end{definition}

\begin{remark}[the pairing is bilinear, not Hermitian]
There is no complex conjugation in $\ip{T}{\varphi}$. It is linear in each slot over $\C$. In
particular $\ip{T}{\varphi}$ is not an inner product, and statements such as ``$\F$ is unitary'' have
no meaning in this framework. The absence of conjugation is what fixes the position of the factors
$2\pi$ in Chapter~\ref{ch:tespered}.
\end{remark}

\begin{proposition}[local estimate]\label{prop:localest}
A linear $T:\D\to\C$ is a distribution if and only if for every compact $K$ there exist $C_{K}<\infty$
and $m_{K}\in\mathbb{N}$ such that
\begin{equation}
|\ip{T}{\varphi}|\le C_{K}\max_{0\le k\le m_{K}}\|\varphi^{(k)}\|_{\infty}
\qquad\text{whenever }\supp\varphi\subset K .
\label{eq:localest}
\end{equation}
\end{proposition}

\begin{proof}[Sketch]
Sufficiency: \eqref{eq:localest} makes continuity immediate, since convergence in $\D$ forces every
derivative to converge uniformly on $K$. Necessity: if no such estimate existed, one could construct
$\varphi_{j}\to0$ in $\D$ with $|\ip{T}{\varphi_{j}}|\ge1$, contradicting continuity.
\end{proof}

The estimate expresses the only kind of control required: on a fixed region, the response is bounded
by finitely many derivatives of the probe. The integer $m_{K}$ is the \emph{order} of $T$ on $K$.

\section{Ordinary functions as distributions}

\begin{theorem}[embedding]\label{thm:embedding}
Let $f\in\Lonloc$. Then
\begin{equation}
\ip{T_{f}}{\varphi}:=\int_{\R}f(x)\varphi(x)\dd x
\end{equation}
defines a distribution. Moreover $T_{f}=T_{g}$ if and only if $f=g$ almost everywhere.
\end{theorem}

\begin{proof}
Since $\varphi$ is supported in a compact $K$, the integral is over $K$ and is finite by local
integrability. The estimate
\begin{equation}
|\ip{T_{f}}{\varphi}|\le\|f\|_{L^{1}(K)}\|\varphi\|_{\infty}
\end{equation}
is \eqref{eq:localest} with $m_{K}=0$. For the last statement, if $\int(f-g)\varphi=0$ for all
$\varphi\in\D$ supported in $K$, then $f-g=0$ almost everywhere on $K$; this is the standard
duality between locally integrable functions and test functions. Taking a sequence of compact sets
exhausting $\R$ gives the result.
\end{proof}

\begin{remark}[values at a point are invisible]
The theorem says that a distribution sees only the almost-everywhere class of an $\Lonloc$ function.
In particular the value $H(0)$ is irrelevant: changing it changes the function on a null set, hence
not the distribution. This is not vagueness; it is a precise statement, and it removes a question
that would otherwise be unanswerable.
\end{remark}

\section{\texorpdfstring{$\delta$}{delta} is not a function}

\begin{theorem}\label{thm:delta-not-L1}
There is no $g\in\Lonloc$ with $T_{g}=\delta$.
\end{theorem}

\begin{proof}
Suppose $T_{g}=\delta$. Let $\varphi\in C_{c}^{\infty}$ be supported in $\R\setminus\{0\}$. Then
$\ip{\delta}{\varphi}=0$, so $\int g\varphi=0$. As $\varphi$ ranges over test functions supported in
$\R\setminus\{0\}$, Theorem~\ref{thm:embedding} gives $g=0$ almost everywhere on $\R\setminus\{0\}$.
Since $\{0\}$ is a null set, $g=0$ almost everywhere on $\R$, hence $T_{g}=0\neq\delta$, a
contradiction.
\end{proof}

So $\delta$ is a genuinely new object: it is not a function with a strange value at one point, it is
not any function at all. Two remarks prevent over-reading this.

\begin{enumerate}
\item ``Not a function'' does not mean ``not defined''. $\delta$ is defined by
\eqref{eq:deltadef}, and it is also a measure: the Dirac point mass, assigning $1$ to a set
containing $0$ and $0$ otherwise. Measures embed into distributions, and $\delta$ is the simplest
member of the intersection.
\item The failure in Theorem~\ref{thm:delta-not-L1} is a failure of \emph{representation}, not of
existence. Every $\Lonloc$ function gives a distribution; $\delta$ is a distribution that no
$\Lonloc$ function gives.
\end{enumerate}

\section{Operations}

\subsection*{Differentiation}

\begin{definition}
$\ip{T'}{\varphi}:=-\ip{T}{\varphi'}$, and $\ip{T^{(n)}}{\varphi}:=(-1)^{n}\ip{T}{\varphi^{(n)}}$.
\end{definition}

This is well defined because $\varphi\in\D$ implies $\varphi'\in\D$, and because
$\varphi_{j}\to\varphi$ in $\D$ implies $\varphi_{j}'\to\varphi'$ in $\D$.

\begin{proposition}[differentiation preserves convergence]\label{prop:derivcont}
If $T_{j}\to T$ in $\D'$, then $T_{j}'\to T'$ in $\D'$.
\end{proposition}

\begin{proof}
$\ip{T_{j}'}{\varphi}=-\ip{T_{j}}{\varphi'}\to-\ip{T}{\varphi'}=\ip{T'}{\varphi}$ for every
$\varphi$.
\end{proof}

\subsection*{Multiplication by a smooth function}

For $a\in C^{\infty}(\R)$ define
\begin{equation}
\ip{aT}{\varphi}:=\ip{T}{a\varphi},
\end{equation}
which is well defined because $a\varphi\in\D$. The operation is associative and compatible with
differentiation.

\begin{proposition}\label{prop:mult-rules}
For $a\in C^{\infty}$ and any distribution $T$,
\begin{align}
a\,\delta&=a(0)\,\delta, \label{eq:adelta}\\
x\,\delta&=0, \label{eq:xdelta}\\
x\,\delta'&=-\delta, \label{eq:xdelta'}\\
x\,\delta^{(n)}&=-n\,\delta^{(n-1)}, \label{eq:xdeltaN}\\
(aT)'&=a'T+aT' . \label{eq:leibniz}
\end{align}
\end{proposition}

\begin{proof}
For \eqref{eq:adelta}: $\ip{a\delta}{\varphi}=\ip{\delta}{a\varphi}=a(0)\varphi(0)=a(0)\ip{\delta}{\varphi}$.

For \eqref{eq:xdelta'}: $\ip{x\delta'}{\varphi}=\ip{\delta'}{x\varphi}=-(x\varphi)'(0)=-\varphi(0)$.

For \eqref{eq:xdeltaN}: since $(x\varphi)^{(n)}=x\varphi^{(n)}+n\varphi^{(n-1)}$, evaluating at $0$
leaves $n\varphi^{(n-1)}(0)$, so
\begin{equation}
\ip{x\delta^{(n)}}{\varphi}=(-1)^{n}n\varphi^{(n-1)}(0)
=-n\,(-1)^{n-1}\varphi^{(n-1)}(0)=-n\ip{\delta^{(n-1)}}{\varphi}.
\end{equation}

For \eqref{eq:leibniz}: using $(a\varphi)'=a'\varphi+a\varphi'$,
\begin{equation}
\ip{(aT)'}{\varphi}=-\ip{T}{(a\varphi)'}
=-\ip{T}{a'\varphi}-\ip{T}{a\varphi'}
=\ip{a'T}{\varphi}+\ip{aT'}{\varphi}. \qedhere
\end{equation}
\end{proof}

\begin{remark}
Formula \eqref{eq:xdeltaN} explains Proposition~\ref{prop:keromega} a second time: multiplying by
$\omega$ lowers the order of $\delta^{(n)}$ by one, so among the $\delta^{(n)}$ only $\delta$ itself
is annihilated.
\end{remark}

\subsection*{Support}

Definition~\ref{def:support} applies verbatim to distributions, with $\supp\varphi\subset U$. The
support of $\delta^{(n)}$ is $\{0\}$ (Proposition~\ref{prop:suppdelta}); the support of $H$ is
$[0,\infty)$; the support of a polynomial is all of $\R$ as a function, but its \emph{spectrum} is
supported at the origin, and the two uses of the word must not be confused.

\section{Translation, scaling, composition}

\begin{proposition}\label{prop:affine}
Let $a\neq0$ and $b\in\R$. Then
\begin{equation}
\ip{\delta_{b}}{\varphi}=\varphi(b),
\qquad
\delta(ax)=\frac{1}{|a|}\,\delta(x),
\qquad
\delta(x-b)=\delta_{b}(x).
\label{eq:delta-scale}
\end{equation}
\end{proposition}

\begin{proof}
For the scaling, test against $\varphi$ and substitute $u=ax$:
\begin{equation}
\ip{\delta(a\cdot)}{\varphi}=\left\langle\delta,\frac{1}{|a|}\varphi\!\left(\frac{\cdot}{a}\right)\right\rangle
=\frac{1}{|a|}\varphi(0)=\left\langle\frac{1}{|a|}\delta,\varphi\right\rangle . \qedhere
\end{equation}
\end{proof}

The factor $|a|^{-1}$ is a Jacobian, not an algebraic reciprocal; the absolute value matters when
$a<0$.

\begin{proposition}[composition]\label{prop:composition}
Let $g\in C^{\infty}$ have only simple zeros $x_{j}$, locally finite in number. Then
\begin{equation}
\delta\big(g(x)\big)=\sum_{j}\frac{\delta(x-x_{j})}{|g'(x_{j})|} .
\label{eq:compose}
\end{equation}
\end{proposition}

\begin{proof}[Sketch]
Near a simple zero $x_{j}$, write $g(x)=(x-x_{j})h_{j}(x)$ with $h_{j}(x_{j})=g'(x_{j})\neq0$; the
change of variable $u=g(x)$ on a small interval around $x_{j}$ has Jacobian
$1/|g'(x_{j})|$, and outside the union of these intervals the test function is annihilated because
$g$ does not vanish there. Summing the contributions gives \eqref{eq:compose}.
\end{proof}

\begin{remark}
The hypothesis of simple zeros cannot be dropped. For $g(x)=x^{2}$ the derivative vanishes at the
zero, the formula divides by zero, and indeed $\delta(x^{2})$ is not of this form.
\end{remark}

\section{Operations that are not available}

\begin{enumerate}
\item \textbf{Products of distributions.} There is no general definition of $ST$ for
$S,T\in\D'$. In particular $\delta^{2}$ has no canonical meaning, and expressions such as
$\|\delta\|^{2}$ or ``the energy of $\delta$'' are meaningless.
\item \textbf{Multiplication by nonsmooth functions.} $H\delta$ has no intrinsic definition: $H$ is
not smooth, and different regularizations of the product give different answers. Choosing $H(0)$
does not resolve this.
\item \textbf{Convolutions of arbitrary distributions.} A compactly supported distribution may be
convolved with any distribution; two arbitrary distributions need not admit a convolution.
\item \textbf{Confusion with measures.} Distributions are not signed measures. $\delta$ is both, but
$\delta'$ is not a measure, and neither is $\PV\frac{1}{\omega}$ in the sense of positive or signed
measure with finite local variation (Proposition~\ref{prop:Kprime}(iv) makes this visible: the
total variation diverges).
\end{enumerate}

\section{Regularization: every distribution is a limit of smooth functions}

One structural fact remains, and it explains why the method of \emph{approximating by ordinary
functions} never loses information.

Let $\rho\in C_{c}^{\infty}$ with $\rho\ge0$, $\int\rho=1$, and set
$\rho_{\ep}(x)=\ep^{-1}\rho(x/\ep)$.

\begin{theorem}\label{thm:regularization}
Let $T$ be a distribution with compact support. For $\ep>0$ define
\begin{equation}
T_{\ep}(x):=\big\langle T,\;\rho_{\ep}(x-\cdot)\big\rangle .
\label{eq:Teps}
\end{equation}
Then $T_{\ep}\in C^{\infty}(\R)$, and $T_{\ep}\to T$ in $\D'(\R)$ as $\ep\downarrow0$.
\end{theorem}

\begin{proof}
\emph{Smoothness.} Let $K$ be a compact set containing $\supp T$ in its interior. For $x$ in a
bounded set, the map $y\mapsto\rho_{\ep}(x-y)$ is supported in a fixed compact neighbourhood of $K$, so
the pairing is legitimate. Differentiation under the pairing is allowed because
$\partial_{x}^{\alpha}\rho_{\ep}(x-y)=(-\partial_{y})^{\alpha}\rho_{\ep}(x-y)$ is again a test
function, and $T$ is continuous; hence $T_{\ep}\in C^{\infty}$.

\emph{Convergence.} Let $\varphi\in C_{c}^{\infty}$. Interchanging the pairing with the integration in
$x$ (a finite sum after expanding on a compact set),
\begin{equation}
\ip{T_{\ep}}{\varphi}
=\int_{\R}\big\langle T,\rho_{\ep}(x-\cdot)\big\rangle\varphi(x)\dd x
=\Big\langle T,\;\int_{\R}\rho_{\ep}(x-\cdot)\varphi(x)\dd x\Big\rangle
=\big\langle T,\;\rho_{\ep}*\varphi\big\rangle ,
\end{equation}
where the last step uses the symmetry $\rho_{\ep}(-y)=\rho_{\ep}(y)$. By
Theorem~\ref{thm:conv}, $\rho_{\ep}*\varphi\to\varphi$ in $C^{\infty}$ uniformly on compact sets, hence
in $\D$, and continuity of $T$ gives $\ip{T_{\ep}}{\varphi}\to\ip{T}{\varphi}$.
\end{proof}

\begin{remark}
The hypothesis of compact support is what allows the mollification to be tested against $T$ without
worrying about behaviour at infinity. For general $T$ a localization argument---multiply by a smooth
cutoff equal to one on a large set---reduces to this case, at the cost of a statement that holds on
compact subsets only, which is exactly what convergence in $\D'$ means.
\end{remark}

\begin{corollary}\label{cor:approx}
Every result in this book that was obtained by regularizing with a specific family
($e^{-\ep|t|}$, $H e^{-\ep t}$, $g_{\sigma}$, $w_{T}$, $K_{\ep}$) is an instance of a single
mechanism: convolving, or pairing, against a mollifier.
\end{corollary}

\section{Summary of the chapter}

\begin{center}
\small
\begin{tabular}{@{}p{3.4cm}p{9.4cm}@{}}
\toprule
Notion & Content \\
\midrule
Test functions & $C_{c}^{\infty}$: smooth (to move derivatives) and compactly supported (to kill
boundary terms) \\
Distribution & continuous linear functional on $\D$; continuity estimated by
\eqref{eq:localest} \\
Embedding & $\Lonloc\hookrightarrow\D'$, injective modulo null sets
(Theorem~\ref{thm:embedding}) \\
$\delta$ & a distribution, not an $\Lonloc$ function
(Theorem~\ref{thm:delta-not-L1}); also a measure \\
Derivative & transfer to the probe with $(-1)^{n}$; continuous for distributional convergence \\
Multiplication & only by smooth functions; \eqref{eq:adelta}--\eqref{eq:leibniz} \\
Forbidden & products of distributions, $\delta^{2}$, $H\delta$, unrestricted convolution \\
\bottomrule
\end{tabular}
\end{center}

\chapter{Tempered Distributions and the Generalized Transform}
\label{ch:tespered}

\begin{quote}
\itshape
The transform cannot be transplanted to $\D'$ because it does not preserve $\D$. It is preserved by
a smaller but more regular space, the Schwartz space, and on the dual of that space the transform is
defined by moving it to the probe. Inversion is proved independently by Gaussian regularization, so
nothing circular enters. The three results of the preceding chapters then become theorems.
\end{quote}

\section{Why the naive definition fails}

The natural attempt is
\begin{equation}
\ip{\F T}{\varphi}:=\ip{T}{\F\varphi},
\label{eq:naive}
\end{equation}
which requires $\F\varphi$ to be a test function whenever $\varphi$ is. It is not.

\begin{proposition}\label{prop:FnotD}
If $\varphi\in C_{c}^{\infty}$ and $\varphi\not\equiv0$, then $\F\varphi\notin C_{c}^{\infty}$.
\end{proposition}

\begin{proof}
For complex $\zeta$ the integral $\ph(\zeta)=\int\varphi(t)e^{-i\zeta t}\dd t$ converges absolutely
because $\varphi$ is compactly supported, and it defines an entire function of $\zeta$. An entire
function vanishing on a nonempty open set vanishes identically, whence $\varphi\equiv0$ by inversion.
So $\ph$ cannot have compact support unless $\varphi\equiv0$.
\end{proof}

\begin{example}
The rectangular window $w_{T}=\mathbf{1}_{[-T,T]}$ has transform $2\sin(T\omega)/\omega$, which
oscillates along the whole real line. Smoothing the corners does not shorten its tail; by
Proposition~\ref{prop:FnotD} no nonzero smooth compactly supported $\varphi$ has a compactly supported
transform.
\end{example}

The conclusion is structural: a space of probes on which $\F$ acts must be invariant under the
transform, and $C_{c}^{\infty}$ is not.

\section{The Schwartz space}

\begin{definition}\label{def:S}
$\Sw(\R)$ is the set of $\varphi\in C^{\infty}(\R)$ such that for all $m,n\in\mathbb{N}$,
\begin{equation}
p_{m,n}(\varphi):=\sup_{x\in\R}(1+|x|)^{m}\big|\varphi^{(n)}(x)\big|<\infty .
\label{eq:seminorm}
\end{equation}
The $p_{m,n}$ are seminorms, and $\varphi_{j}\to\varphi$ in $\Sw$ means
$p_{m,n}(\varphi_{j}-\varphi)\to0$ for all $m,n$.
\end{definition}

Equivalently: every derivative decays faster than any negative power of $|x|$.

\begin{example}\label{ex:S}
$e^{-x^{2}}\in\Sw$; every element of $C_{c}^{\infty}$ belongs to $\Sw$; but $1$ and $x^{n}$ do not
(they do not decay), $e^{x^{2}}$ does not (it grows), and $1/(1+x^{2})$ does not, although it is
integrable: it decays only like $x^{-2}$, and its derivatives decay more slowly still. Integrability
and rapid decay are different requirements.
\end{example}

\section{The transform acts on the Schwartz space}

Two identities carry the whole proof.

\begin{lemma}\label{lem:two-identities}
For $\varphi\in\Sw$ and $n,m\in\mathbb{N}$,
\begin{equation}
\partial_{\omega}^{\,n}\ph(\omega)=\F\big[(-it)^{n}\varphi\big](\omega),
\qquad
\omega^{m}\,\ph(\omega)=i^{-m}\,\F\big[\varphi^{(m)}\big](\omega).
\label{eq:two-identities}
\end{equation}
\end{lemma}

\begin{proof}
The first is Lemma~\ref{lem:diffint} iterated. For the second, Theorem~\ref{thm:derivt} gives
$\F[\varphi^{(m)}]=(i\omega)^{m}\ph$, and multiplying by $i^{-m}$ yields \eqref{eq:two-identities}.
\end{proof}

\begin{theorem}\label{prop:F-S}
$\F$ maps $\Sw$ continuously into $\Sw$; that is, $\ph\in\Sw$ for $\varphi\in\Sw$, and for all
$m,n$ there are $C$ and finitely many indices $(m_{k},n_{k})$ with
$p_{m,n}(\ph)\le C\sum p_{m_{k},n_{k}}(\varphi)$.
\end{theorem}

\begin{proof}
Combining the two identities,
\begin{equation}
\omega^{m}\partial_{\omega}^{\,n}\ph(\omega)
=i^{-m}\int_{\R}\partial_{t}^{m}\big(t^{n}\varphi(t)\big)e^{-i\omega t}\dd t ,
\end{equation}
hence
\begin{equation}
\big|\omega^{m}\partial_{\omega}^{\,n}\ph(\omega)\big|
\le\big\|\partial_{t}^{m}(t^{n}\varphi)\big\|_{1} .
\label{eq:Fbound}
\end{equation}
The right side does not depend on $\omega$. Since $\varphi\in\Sw$, the function $t^{n}\varphi$ lies in
$\Sw$, so each of its derivatives is integrable, and the right side of \eqref{eq:Fbound} is finite.
Taking the supremum over $\omega$ and using $(1+|\omega|)^{m}\le C_{m}\sum_{k\le m+1}|\omega|^{k}$
yields the stated bound.
\end{proof}

The proof shows why $\Sw$ is the right space. The two operations that appear in the transform---taking
a derivative and multiplying by the variable---exchange places under $\F$, and $\Sw$ is closed under
both.

\section{Inversion}

The inversion formula cannot be obtained from
$\int e^{i\omega(t-s)}\dd\omega=2\pi\delta(t-s)$ without circularity: that identity is itself a
distributional statement about $\delta$. The following proof uses only the Gaussian integral, Fubini,
and the approximate identities of Chapter~\ref{ch:limits}.

\begin{lemma}[Gaussian integral]\label{lem:gauss}
For $\eta>0$ and $x\in\R$,
\begin{equation}
\int_{\R}e^{-\eta\omega^{2}/2}e^{i\omega x}\dd\omega=\sqrt{\frac{2\pi}{\eta}}\;e^{-x^{2}/(2\eta)} .
\label{eq:gaussint}
\end{equation}
\end{lemma}

\begin{proof}
Complete the square in the exponent and use $\int e^{-u^{2}/2}\dd u=\sqrt{2\pi}$.
\end{proof}

\begin{theorem}[Fourier inversion]\label{thm:inversion}
Let $\varphi\in\Sw$. Then
\begin{equation}
\varphi(t)=\frac{1}{2\pi}\int_{\R}\ph(\omega)e^{i\omega t}\dd\omega
\qquad\text{for every }t\in\R .
\label{eq:inversion}
\end{equation}
\end{theorem}

\begin{proof}
For $\eta>0$ define
\begin{equation}
I_{\eta}(t)=\frac{1}{2\pi}\int_{\R}\ph(\omega)e^{i\omega t}e^{-\eta\omega^{2}/2}\dd\omega .
\end{equation}
Insert $\ph(\omega)=\int\varphi(s)e^{-i\omega s}\dd s$. The double integral is absolutely convergent,
because the modulus of the integrand is $|\varphi(s)|e^{-\eta\omega^{2}/2}$ and $\varphi\in\Lone$,
while the Gaussian is integrable in $\omega$. Fubini therefore applies and
\begin{align}
I_{\eta}(t)
&=\frac{1}{2\pi}\int_{\R}\varphi(s)\left[\int_{\R}e^{-\eta\omega^{2}/2}e^{i\omega(t-s)}\dd\omega
\right]\dd s \notag\\
&=\frac{1}{2\pi}\int_{\R}\varphi(s)\sqrt{\frac{2\pi}{\eta}}\,e^{-(t-s)^{2}/(2\eta)}\dd s
=\big(\rho_{\eta}*\varphi\big)(t),
\label{eq:conv}
\end{align}
where $\rho_{\eta}(u)=(2\pi\eta)^{-1/2}e^{-u^{2}/(2\eta)}$ is the Gaussian approximate identity.

Two limits are now taken. On the left, since $\ph\in\Sw\subset\Lone$ and
$|e^{-\eta\omega^{2}/2}|\le1$, dominated convergence with the $\eta$-independent bound $|\ph|$ gives
\begin{equation}
I_{\eta}(t)\longrightarrow\frac{1}{2\pi}\int_{\R}\ph(\omega)e^{i\omega t}\dd\omega .
\end{equation}
On the right, Theorem~\ref{thm:conv} gives $(\rho_{\eta}*\varphi)(t)\to\varphi(t)$. Equating the two
limits yields \eqref{eq:inversion}.
\end{proof}

\begin{theorem}[two transforms give a reflection]\label{thm:F2}
For $\varphi\in\Sw$,
\begin{equation}
\F^{2}\varphi(t)=2\pi\,\varphi(-t).
\label{eq:F2}
\end{equation}
\end{theorem}

\begin{proof}
By definition $\F^{2}\varphi(t)=\int_{\R}\ph(\omega)e^{-i\omega t}\dd\omega$, and
\eqref{eq:inversion} with $t$ replaced by $-t$ identifies this with $2\pi\varphi(-t)$.
\end{proof}

\begin{remark}
The order of the argument matters. Inversion is proved first, without reference to $\delta$; only
afterwards may one use $\F^{2}=2\pi\,(\text{reflection})$ to obtain statements about $\delta$ cheaply.
Using the distributional identity $\int e^{i\omega t}\dd\omega=2\pi\delta(t)$ to prove inversion would
reverse the implication.
\end{remark}

\section{Tempered distributions and the dual definition}

\begin{definition}\label{def:Sprime}
A \emph{tempered distribution} is a continuous linear functional on $\Sw(\R)$. The space is denoted
$\Sw'(\R)$. Continuity is equivalent to the existence of $C<\infty$ and finitely many indices such
that
\begin{equation}
|\ip{T}{\varphi}|\le C\sum_{m\le M,\ n\le N}p_{m,n}(\varphi)
\qquad\text{for all }\varphi\in\Sw .
\label{eq:Sprime-est}
\end{equation}
\end{definition}

\begin{example}\label{ex:tempered}
Bounded functions are tempered, since $|\ip{f}{\varphi}|\le\|f\|_{\infty}\|\varphi\|_{1}$ and
$\|\varphi\|_{1}$ is controlled by $p_{2,0}$. Functions of polynomial growth are tempered, with the
index $N$ equal to the growth order plus $2$. $\delta^{(n)}$ is tempered, since
$|\ip{\delta^{(n)}}{\varphi}|=|\varphi^{(n)}(0)|\le p_{0,n}(\varphi)$. The principal value is
tempered, the estimate near the origin using one derivative of $\varphi$ and the estimate at
infinity using rapid decay.
\end{example}

\begin{proposition}[not every distribution is tempered]\label{prop:ex2}
The function $e^{x^{2}}$ is locally integrable, hence defines a distribution, but it is not a
tempered distribution.
\end{proposition}

\begin{proof}
Take $0\not\equiv\psi\in C_{c}^{\infty}$ nonnegative and set $\psi_{R}(x)=\psi(x-R)$. The seminorms
$p_{m,n}(\psi_{R})$ grow at most polynomially in $R$, because translation only multiplies the weight
$(1+|x|)^{m}$ by a factor $O(R^{m})$. On the other hand
\begin{equation}
\ip{e^{x^{2}}}{\psi_{R}}=\int e^{x^{2}}\psi(x-R)\dd x\ge e^{(R-r)^{2}}\int\psi ,
\end{equation}
where $r$ is the radius of $\supp\psi$, which grows faster than any polynomial. No estimate of the
form \eqref{eq:Sprime-est} can hold.
\end{proof}

\begin{definition}[generalized Fourier transform]\label{def:FTdist}
For $T\in\Sw'$ define $\F T\in\Sw'$ by
\begin{equation}
\ip{\F T}{\varphi}:=\ip{T}{\F\varphi},
\qquad \varphi\in\Sw .
\label{eq:dualdef}
\end{equation}
\end{definition}

\begin{proposition}[consistency with the classical transform]\label{prop:dual-consistency}
If $f,\varphi\in\Sw$, then $\ip{\F f}{\varphi}=\ip{f}{\F\varphi}$; consequently \eqref{eq:dualdef}
extends the classical transform.
\end{proposition}

\begin{proof}
By Fubini, justified by absolute convergence,
\begin{equation}
\ip{\F f}{\varphi}
=\iint f(t)e^{-i\omega t}\varphi(\omega)\dd t\dd\omega
=\int_{\R}f(t)\Big[\int_{\R}\varphi(\omega)e^{-i\omega t}\dd\omega\Big]\dd t
=\int_{\R}f(t)\ph(t)\dd t=\ip{f}{\F\varphi} . \qedhere
\end{equation}
\end{proof}

\begin{remark}[three details that must be checked]
First, \eqref{eq:dualdef} is well posed because $\varphi\in\Sw$ implies
$\F\varphi\in\Sw$ (Theorem~\ref{prop:F-S}). Second, the right side defines a continuous linear
functional of $\varphi$, hence an element of $\Sw'$. Third, there is no factor $1/(2\pi)$ and no
conjugate on the right; a formula with either present belongs to a different convention.
\end{remark}

\begin{theorem}[inverse transform]\label{thm:inverse}
$\F:\Sw'\to\Sw'$ is a linear isomorphism, with inverse
\begin{equation}
\F^{-1}T=\frac{1}{2\pi}\,\F\big(T(-\cdot)\big).
\end{equation}
\end{theorem}

\begin{proof}
By Theorem~\ref{thm:F2} in dual form, $\F^{2}T=2\pi\,T(-\cdot)$; solving gives the stated inverse.
\end{proof}

\section{Continuity, and the gaps it closes}

\begin{theorem}[the transform is continuous]\label{thm:F-cont}
If $T_{j}\to T$ in $\Sw'$, then $\F T_{j}\to\F T$ in $\Sw'$.
\end{theorem}

\begin{proof}
For fixed $\varphi\in\Sw$,
$\ip{\F T_{j}}{\varphi}=\ip{T_{j}}{\F\varphi}\to\ip{T}{\F\varphi}=\ip{\F T}{\varphi}$.
\end{proof}

This is the step missing from Chapters~\ref{ch:constant} and \ref{ch:step}, where a time-domain limit
and a frequency-domain limit were computed separately without a bridge between them.

\begin{example}[completing the constant]
For every $\varphi\in\Sw$, $|e^{-\ep|t|}\varphi(t)|\le|\varphi(t)|\in\Lone$ and
$e^{-\ep|t|}\varphi(t)\to\varphi(t)$ pointwise, so $e^{-\ep|\cdot|}\to1$ in $\Sw'$ by dominated
convergence. By Theorems~\ref{thm:F-cont} and \ref{prop:F-S},
\begin{equation}
\F\big[e^{-\ep|\cdot|}\big]\longrightarrow\F[1]
\qquad\text{in }\Sw',
\end{equation}
while Theorem~\ref{thm:exptransform} and Theorem~\ref{thm:sampling} give
$\F[e^{-\ep|\cdot|}]=G_{\ep}\to2\pi\delta$. Uniqueness of limits in $\Sw'$ yields
$\F[1]=2\pi\delta$. Compare \eqref{eq:F1}: the same statement, now proved rather than posited.
\end{example}

\section{The three results}

\begin{theorem}\label{thm:main}
In $\Sw'$:
\begin{align}
\F[\delta]&=1, \qquad \F[1]=2\pi\delta, \label{eq:Fdelta}\\
\F[t^{\,n}]&=2\pi i^{\,n}\delta^{(n)},\qquad n=0,1,2,\dots, \label{eq:Ftn}\\
\F[H]&=\pi\delta+\PV\frac{1}{i\omega}. \label{eq:FHgen}
\end{align}
\end{theorem}

\begin{proof}
\emph{(i) $\delta$ and $1$.} For $\varphi\in\Sw$,
\begin{equation}
\ip{\F\delta}{\varphi}=\ip{\delta}{\ph}=\ph(0)=\int_{\R}\varphi(t)\dd t=\ip{1}{\varphi},
\end{equation}
so $\F\delta=1$. For the second identity,
\begin{equation}
\ip{\F1}{\varphi}=\ip{1}{\ph}=\int_{\R}\ph(t)\dd t=2\pi\varphi(0)=\ip{2\pi\delta}{\varphi},
\end{equation}
where the middle equality is \eqref{eq:inversion} at $t=0$. This is the only place where inversion is
used, and it is the reason inversion had to be proved first.

\emph{(ii) The two differentiation rules.} First, from $\F[\varphi']=i\omega\ph$ evaluated at
$\omega=t$ one gets $(\varphi')^{\wedge}(t)=it\,\ph(t)$, that is
$t\,\ph(t)=-i(\varphi')^{\wedge}(t)$. Hence
\begin{equation}
\ip{\F(tT)}{\varphi}=\ip{T}{t\ph}
=-i\ip{T}{(\varphi')^{\wedge}}
=-i\ip{\widehat T}{\varphi'}
=i\ip{\widehat T'}{\varphi},
\end{equation}
so
\begin{equation}
\F(tT)=i\,\partial_{\omega}\widehat T .
\label{eq:multt-dist}
\end{equation}
Second, from $\ph'(\omega)=-i(t\varphi)^{\wedge}(\omega)$,
\begin{equation}
\ip{\F(T')}{\varphi}=\ip{T'}{\ph}=-\ip{T}{\ph'}
=i\ip{\widehat T}{t\varphi}
=i\ip{t\widehat T}{\varphi},
\end{equation}
so
\begin{equation}
\F(T')=i\omega\,\widehat T .
\label{eq:deriv-dist}
\end{equation}
Iterating \eqref{eq:multt-dist} on $\F[1]=2\pi\delta$ yields \eqref{eq:Ftn}.

\emph{(iii) The step.} $H_{\ep}=H\,e^{-\ep t}\to H$ in $\Sw'$ because
$|H_{\ep}\varphi|\le|\varphi|\in\Lone$; and by Theorem~\ref{thm:Heps} together with
Theorems~\ref{thm:realpart} and \ref{thm:imagpart},
$\F[H_{\ep}]\to\pi\delta-i\PV\frac{1}{\omega}$. Theorem~\ref{thm:F-cont} identifies the limit of the
transforms with $\F[H]$, and uniqueness gives \eqref{eq:FHgen}.
\end{proof}

\section{Other regularizations}

The theory must be independent of the chosen regularizing family. Two checks.

\subsection*{Gaussian}

For $g_{\sigma}(t)=e^{-\sigma t^{2}/2}$ one has, by the differential-equation method or
Lemma~\ref{lem:gauss},
\begin{equation}
\widehat{g_{\sigma}}(\omega)=\sqrt{\frac{2\pi}{\sigma}}\,e^{-\omega^{2}/(2\sigma)} .
\end{equation}
Its total mass is
$\sqrt{2\pi/\sigma}\cdot\sqrt{2\pi\sigma}=2\pi$, and it has width of order $\sqrt{\sigma}$ and height
of order $\sigma^{-1/2}$; hence $\widehat{g_{\sigma}}\to2\pi\delta$. The Gaussian route is entirely
independent of the exponential route of Chapter~\ref{ch:constant} and gives the same limit.

\subsection*{Rectangular windows, and sinc}

For $w_{T}=\mathbf{1}_{[-T,T]}$ one has $\widehat{w_{T}}(\omega)=2\sin(T\omega)/\omega$. Since
$w_{T}\to1$ in $\Sw'$ (dominated convergence with bound $|\varphi|$), Theorem~\ref{thm:F-cont} gives
\begin{equation}
\widehat{w_{T}}\longrightarrow\F[1]=2\pi\delta
\qquad\text{in }\Sw'.
\end{equation}

\begin{remark}[why Chapter~\ref{ch:limits} does not apply]
The sinc kernels are not nonnegative, they change sign infinitely often, and
$\int_{\R}\big|2\sin(T\omega)/\omega\big|\dd\omega$ grows like $\log T$. Neither the hypotheses of
Theorem~\ref{thm:sampling-general} nor those of Proposition~\ref{prop:general-kernels} hold. Their
convergence comes instead from the continuity of $\F$ applied to the convergent family $w_{T}\to1$.
There is no contradiction with the failure of pointwise convergence noted in
Proposition~\ref{prop:nolimit}: that failure concerns the pointwise behaviour of the truncated
integrals, this concerns distributional convergence of the transforms.
\end{remark}

\begin{remark}[what the extension does and does not do]
It does not assign values to divergent integrals. It does the following: it reinterprets a signal as
a functional, defines $\F$ on functionals by duality, proves that the definition is continuous and
invertible, and then observes that the three identities hold simultaneously. The extension is
conservative---it agrees with the classical transform wherever the latter exists---and it is
canonical, because inversion together with the pairing leaves no freedom in the definition.
\end{remark}

\section{Plancherel, and what the delta does not have}

The generalized transform extends the classical one. It does not extend Plancherel's identity to
$\delta$, and it is worth being precise about why.

\begin{theorem}[Plancherel]\label{thm:plancherel}
For $f,g\in\Sw$,
\begin{equation}
\int_{\R}f(t)\overline{g(t)}\dd t
=\frac{1}{2\pi}\int_{\R}\fh(\omega)\overline{\gh(\omega)}\dd\omega .
\label{eq:plancherel-pair}
\end{equation}
In particular $\|f\|_{2}^{2}=(2\pi)^{-1}\|\fh\|_{2}^{2}$, and $\F$ extends uniquely to a bounded
operator $L^{2}(\R)\to L^{2}(\R)$ with $\|\F f\|_{2}=\sqrt{2\pi}\,\|f\|_{2}$.
\end{theorem}

\begin{proof}
Since $g=\frac{1}{2\pi}\int\gh(\omega)e^{i\omega\cdot}\dd\omega$ by inversion, and since
$\gh\in\Sw\subset\Lone$, Fubini gives
\begin{align}
\int_{\R}f(t)\overline{g(t)}\dd t
&=\int_{\R}f(t)\overline{\frac{1}{2\pi}\int_{\R}\gh(\omega)e^{i\omega t}\dd\omega}\dd t \notag\\
&=\frac{1}{2\pi}\int_{\R}\overline{\gh(\omega)}\Big[\int_{\R}f(t)e^{-i\omega t}\dd t\Big]\dd\omega
=\frac{1}{2\pi}\int_{\R}\fh(\omega)\overline{\gh(\omega)}\dd\omega .
\end{align}
Taking $g=f$ gives the norm identity. The extension: $\Sw$ is dense in $L^{2}$, and on $\Sw$ the map
$\F/\sqrt{2\pi}$ preserves the norm, so it extends uniquely by continuity to an isometry of $L^{2}$;
invertibility follows from \eqref{eq:F2}.
\end{proof}

\begin{corollary}\label{cor:no-energy}
The constant $1$ and the distribution $\delta$ are not elements of $L^{2}$, and neither can be
assigned a finite energy. For $f_{\ep}=e^{-\ep|t|}$,
\begin{equation}
\|f_{\ep}\|_{2}^{2}=\frac{1}{\ep}\longrightarrow\infty,
\qquad
\frac{1}{2\pi}\|G_{\ep}\|_{2}^{2}=\frac{1}{\ep},
\end{equation}
so the constant value $2\pi$ of the area of $G_{\ep}$ carries no information about
\eqref{eq:plancherel-pair} in the limit.
\end{corollary}

\begin{proof}
$\int_{\R}1\dd t=\infty$; for $f_{\ep}$ the two sides of
\eqref{eq:plancherel-pair} are computed by direct integration, using
$\int_{\R}(\ep^{2}+\omega^{2})^{-2}\dd\omega=\pi/(2\ep^{3})$.
\end{proof}

\begin{remark}[three layers, and their intersection]
On $L^{1}$ the transform is an absolutely convergent integral; on $L^{2}$ it is an isometry obtained
by density; on $\Sw'$ it is defined by duality and exists for objects---$\delta$, $1$, $t^{n}$,
$H$---that belong to neither $L^{1}$ nor $L^{2}$. The three definitions agree on their common domain,
which is $\Sw$ at least, and this agreement is the content of
Proposition~\ref{prop:dual-consistency}. The constructions do not replace one another; they are the
same operator seen in three topologies.
\end{remark}

\section{Summary of the chapter}

\begin{center}
\small
\begin{tabular}{@{}p{3.4cm}p{9.4cm}@{}}
\toprule
Step & Content \\
\midrule
Obstruction & $\F$ does not preserve $C_{c}^{\infty}$ (Proposition~\ref{prop:FnotD}) \\
Replacement & $\Sw$, closed under differentiation and multiplication (Definition~\ref{def:S}) \\
Closure & $\F:\Sw\to\Sw$ continuous (Theorem~\ref{prop:F-S}) \\
Inversion & proved by Gaussian regularization, with no appeal to $\delta$
(Theorem~\ref{thm:inversion}) \\
Definition & $\ip{\F T}{\varphi}=\ip{T}{\F\varphi}$ on $\Sw'$
(Definition~\ref{def:FTdist}) \\
Continuity & $T_{j}\to T$ implies $\F T_{j}\to\F T$ (Theorem~\ref{thm:F-cont}) \\
Harvest & $\F[\delta]=1$, $\F[1]=2\pi\delta$, $\F[t^{n}]=2\pi i^{n}\delta^{(n)}$,
$\F[H]=\pi\delta+\PV\frac{1}{i\omega}$ (Theorem~\ref{thm:main}) \\
Consistency & Gaussian and sinc families give the same limits \\
\bottomrule
\end{tabular}
\end{center}

\appendix

\chapter{Tables}
\label{app:tables}

All formulas use the convention
$\F[f](\omega)=\int f(t)e^{-i\omega t}\dd t$ with $2\pi$ on the inverse transform, and pairings
without complex conjugation.

\section{Transform pairs}

\begin{center}
\small
\begin{tabular}{@{}lll@{}}
\toprule
Time domain & Frequency domain & Where proved \\
\midrule
$\delta(t)$ & $1$ & \eqref{eq:Fdelta} \\
$\delta^{(n)}(t)$ & $(i\omega)^{n}$ & \eqref{eq:Fdelta} \\
$1$ & $2\pi\delta(\omega)$ & \eqref{eq:Fdelta} \\
$t^{\,n}$ & $2\pi i^{\,n}\delta^{(n)}(\omega)$ & \eqref{eq:Ftn} \\
$\delta(t-a)$ & $e^{-ia\omega}$ & \eqref{eq:timeshift} \\
$e^{i\omega_{0}t}$ & $2\pi\delta(\omega-\omega_{0})$ & \eqref{eq:modulation} \\
$\cos(\omega_{0}t)$ & $\pi\big[\delta(\omega-\omega_{0})+\delta(\omega+\omega_{0})\big]$ & linearity \\
$\sin(\omega_{0}t)$ & $(\pi/i)\big[\delta(\omega-\omega_{0})-\delta(\omega+\omega_{0})\big]$ & linearity \\
$H(t)$ & $\pi\delta(\omega)-i\,\PV(1/\omega)$ & \eqref{eq:FHgen} \\
$\sgn t$ & $-2i\,\PV(1/\omega)$ & \eqref{eq:sgn} \\
$\PV(1/t)$ & $-i\pi\,\sgn\omega$ & duality, \eqref{eq:F2} \\
$e^{-\ep|t|}$ & $2\ep/(\ep^{2}+\omega^{2})$ & Theorem~\ref{thm:exptransform} \\
$H(t)e^{-\ep t}$ & $1/(\ep+i\omega)$ & Theorem~\ref{thm:Heps} \\
$e^{-at^{2}/2}$, $a>0$ & $\sqrt{2\pi/a}\;e^{-\omega^{2}/(2a)}$ & Lemma~\ref{lem:gauss} \\
\bottomrule
\end{tabular}
\end{center}

\begin{remark}[the sine entry]
Write $\sin(\omega_{0}t)=\frac{1}{2i}(e^{i\omega_{0}t}-e^{-i\omega_{0}t})$ and transform term by
term. The coefficient is $\pi$, not $2\pi$, because sine is half the difference of two exponentials.
Some tables list $\pi i[\delta(\omega-\omega_{0})-\delta(\omega+\omega_{0})]$, which corresponds to
the opposite sign convention in the exponent; the two differ by a sign, and neither oddness nor
vanishing at the origin distinguishes them, so the derivation must be redone rather than recalled.
\end{remark}

\section{Operational rules}

\begin{center}
\small
\begin{tabular}{@{}lll@{}}
\toprule
Rule & Formula & Hypothesis \\
\midrule
time shift & $\F[f(\cdot-a)]=e^{-ia\omega}\fh$ & \eqref{eq:timeshift} \\
modulation & $\F[e^{i\omega_{0}t}f]=\fh(\omega-\omega_{0})$ & \eqref{eq:modulation} \\
dilation & $\F[f(a\,\cdot)]=|a|^{-1}\fh(\omega/a)$ & $a\neq0$ \\
multiply by $t$ & $\F[tf]=i\,\partial_{\omega}\fh$ & $f\in\Sw$ \\
differentiate & $\F[f']=i\omega\fh$ & $f\in\Sw$ \\
distributional & $\F(tT)=i\,\partial_{\omega}\widehat T$ & \eqref{eq:multt-dist} \\
distributional & $\F(T')=i\omega\widehat T$ & \eqref{eq:deriv-dist} \\
twice & $\F^{2}T=2\pi\,T(-\cdot)$ & \eqref{eq:F2} \\
inverse & $\F^{-1}T=(2\pi)^{-1}\F(T(-\cdot))$ & Theorem~\ref{thm:inverse} \\
\midrule
scale $\delta$ & $\delta(ax)=|a|^{-1}\delta(x)$ & $a\neq0$ \\
compose $\delta$ & $\delta(g(x))=\sum_{j}\delta(x-x_{j})/|g'(x_{j})|$ & $g$ has simple zeros \\
multiply $\delta$ & $a\,\delta=a(0)\,\delta$ & $a\in C^{\infty}$ \\
multiply by $x$ & $x\,\delta^{(n)}=-n\,\delta^{(n-1)}$ & \eqref{eq:xdeltaN} \\
Leibniz & $(aT)'=a'T+aT'$ & $a\in C^{\infty}$ \\
jumps & $Df=f'_{\mathrm{cl}}+\sum_{j}J_{j}\delta_{a_{j}}$ & $f$ piecewise $C^{1}$ \\
\bottomrule
\end{tabular}
\end{center}

\section{Integrals used}

\begin{center}
\small
\begin{tabular}{@{}lll@{}}
\toprule
Integral & Value & Used in \\
\midrule
$\displaystyle\int_{0}^{\infty}e^{-at}\dd t$ & $1/a$ \ ($\Re a>0$) & Theorem~\ref{thm:exptransform} \\
$\displaystyle\int_{\R}\frac{\dd\omega}{\ep^{2}+\omega^{2}}$ & $\pi/\ep$ & Proposition~\ref{prop:three} \\
$\displaystyle\int_{\R}\frac{\ep\dd\omega}{\ep^{2}+\omega^{2}}$ & $\pi$ & \eqref{eq:masspi} \\
$\displaystyle\int_{\R}\frac{2\ep\dd\omega}{\ep^{2}+\omega^{2}}$ & $2\pi$ & \eqref{eq:three} \\
$\displaystyle\int_{-a}^{a}\frac{\ep\dd\omega}{\pi(\ep^{2}+\omega^{2})}$ &
$\frac{2}{\pi}\arctan\frac{a}{\ep}$ & \eqref{eq:concentration} \\
$\displaystyle\int_{\R}e^{-u^{2}/2}\dd u$ & $\sqrt{2\pi}$ & Lemma~\ref{lem:gauss} \\
$\displaystyle\int_{\R}e^{-\eta\omega^{2}/2}e^{i\omega x}\dd\omega$ &
$\sqrt{2\pi/\eta}\,e^{-x^{2}/(2\eta)}$ & \eqref{eq:gaussint} \\
$\displaystyle\int_{-T}^{T}e^{-i\omega t}\dd t$ & $2\sin(\omega T)/\omega$ & \eqref{eq:IT} \\
$\displaystyle\int_{\R}\frac{2\ep x}{\pi(x^{2}+\ep^{2})^{2}}\dd x$ & $1/(\pi\ep)$ & Proposition~\ref{prop:Kprime} \\
\bottomrule
\end{tabular}
\end{center}

\section{Which space does an object belong to}

\begin{center}
\small
\begin{tabular}{@{}lllll@{}}
\toprule
Object & $\Lone$ & $\Ltwo$ & $\D'$ & $\Sw'$ \\
\midrule
$e^{-\ep|t|}$ & yes & yes & yes & yes \\
$1$ & no & no & yes & yes \\
$t^{\,n}$ & no & no & yes & yes \\
$H$ & no & no & yes & yes \\
$\delta$, $\delta^{(n)}$ & no & no & yes & yes \\
$\PV(1/\omega)$ & no & no & yes & yes \\
$1/(1+t^{2})$ & yes & yes & yes & no \\
$e^{t^{2}}$ & no & no & yes & no \\
\bottomrule
\end{tabular}
\end{center}

\begin{remark}
Two rows carry most of the information. The constant $1$ is in neither $\Lone$ nor $\Ltwo$, which is
why neither the elementary theory nor Plancherel can be applied to it; it is in $\Sw'$, which is why
the generalized transform can. The function $e^{t^{2}}$ is locally integrable and therefore defines a
distribution, but it is not tempered and consequently has no Fourier transform in this framework.
Membership is a property of the space, not of the formula.
\end{remark}

\section{Conventions in use}

\begin{center}
\small
\begin{tabular}{@{}llll@{}}
\toprule
Convention & Forward transform & $\F[1]$ & Relation to this text \\
\midrule
angular (here) & $\displaystyle\int f(t)e^{-i\omega t}\dd t$ & $2\pi\delta(\omega)$ & --- \\
symmetric & $\displaystyle\frac{1}{\sqrt{2\pi}}\int f(t)e^{-i\omega t}\dd t$ &
$\sqrt{2\pi}\,\delta(\omega)$ & multiply by $\frac{1}{\sqrt{2\pi}}$ \\
frequency & $\displaystyle\int f(t)e^{-2\pi i\nu t}\dd t$ & $\delta(\nu)$ & set $\omega=2\pi\nu$ \\
\bottomrule
\end{tabular}
\end{center}

The conversion rule is
\begin{equation}
\F_{\nu}[f](\nu)=\F_{\omega}[f](2\pi\nu),
\qquad
\delta(2\pi\nu)=\frac{1}{2\pi}\delta(\nu).
\end{equation}

\clearpage
\section{Where each statement is proved}

\begin{center}
\small
\begin{tabular}{@{}lll@{}}
\toprule
Statement & First obtained & Fully justified \\
\midrule
$\F[e^{-\ep|\cdot|}]=2\ep/(\ep^{2}+\omega^{2})$ & Theorem~\ref{thm:exptransform} & same (classical) \\
$\int K_{\ep}\varphi\to\varphi(0)$ & Theorem~\ref{thm:sampling} & Theorem~\ref{thm:sampling-general} \\
$\F[1]=2\pi\delta$ & \eqref{eq:F1} & Theorem~\ref{thm:main} \\
$\ip{\delta^{(n)}}{\varphi}=(-1)^{n}\varphi^{(n)}(0)$ & \eqref{eq:deltan} & \eqref{eq:xdeltaN}--\eqref{eq:leibniz} \\
$\F[t^{\,n}]=2\pi i^{\,n}\delta^{(n)}$ & Theorem~\ref{thm:tn} & Theorem~\ref{thm:main} \\
$\F[H]=\pi\delta+\PV(1/i\omega)$ & \eqref{eq:FH} & Theorem~\ref{thm:main} \\
$\F[\sgn t]=-2i\PV(1/\omega)$ & \eqref{eq:sgn} & Theorem~\ref{thm:main} \\
$\F[\PV(1/t)]=-i\pi\sgn\omega$ & Proposition~\ref{prop:PVtransform} & Theorem~\ref{thm:F2} \\
inversion $\varphi=\frac{1}{2\pi}\int\ph e^{i\omega\cdot}$ & Theorem~\ref{thm:inversion} & same \\
$\F:\Sw\to\Sw$ continuous & Theorem~\ref{prop:F-S} & same \\
$\ip{\F T}{\varphi}=\ip{T}{\F\varphi}$ & \eqref{eq:dualdef} & Proposition~\ref{prop:dual-consistency} \\
Plancherel $\|f\|_{2}^{2}=(2\pi)^{-1}\|\fh\|_{2}^{2}$ & Theorem~\ref{thm:plancherel} & same \\
\bottomrule
\end{tabular}
\end{center}

The second column records where a statement is first \emph{computed}, the third where it is
\emph{derived from the definitions}. The gap between them is the content of
Chapters~\ref{ch:limits}--\ref{ch:tespered}: the early chapters compute with rules that are only
justified later, and the later ones supply the justification without changing any of the numbers.

\chapter{Exercises with Solutions}
\label{app:exercises}

Each exercise certifies one specific step rather than a topic. Solutions are complete enough to be
checked line by line.

\begin{exercise}\label{ex1}
Fix $\omega\neq0$ and show that $I_{T}(\omega)=2\sin(\omega T)/\omega$ has no limit as
$T\to\infty$.
\end{exercise}

\begin{proof}[Solution]
Let $a=|\omega|$ and set $T_{n}=(\pi/2+2\pi n)/a$, $S_{n}=(3\pi/2+2\pi n)/a$. Then
$\sin(aT_{n})=1$ and $\sin(aS_{n})=-1$. Since $\sin(\omega T)=\operatorname{sgn}(\omega)\sin(aT)$,
\begin{equation}
I_{T_{n}}(\omega)=\operatorname{sgn}(\omega)\frac{2}{a},
\qquad
I_{S_{n}}(\omega)=-\operatorname{sgn}(\omega)\frac{2}{a},
\end{equation}
and both $T_{n},S_{n}\to\infty$. Two subsequences with distinct limits preclude convergence.
\end{proof}

\begin{exercise}\label{ex2}
Compute $\int_{-a}^{a}G_{\ep}(\omega)\dd\omega$ and explain why a fixed interval eventually contains
almost all the mass.
\end{exercise}

\begin{proof}[Solution]
\begin{equation}
\int_{-a}^{a}\frac{2\ep}{\ep^{2}+\omega^{2}}\dd\omega
=2\left[\arctan\frac{\omega}{\ep}\right]_{-a}^{a}
=4\arctan\frac{a}{\ep}\longrightarrow4\cdot\frac{\pi}{2}=2\pi .
\end{equation}
The order of quantifiers is the point: $a$ is fixed and $\ep\downarrow0$, so $\arctan(a/\ep)\to\pi/2$.
Sending $a\to\infty$ first would instead compute the total mass of each fixed member, which is a
different statement.
\end{proof}

\begin{exercise}\label{ex3}
For $\varphi(x)=(1+x)e^{-x^{2}}$ compute $\ip{\delta}{\varphi}$, $\ip{\delta'}{\varphi}$,
$\ip{\delta''}{\varphi}$.
\end{exercise}

\begin{proof}[Solution]
$\varphi(0)=1$. Differentiating,
$\varphi'(x)=e^{-x^{2}}\big(1-2x-2x^{2}\big)$, so $\varphi'(0)=1$. A second differentiation at $0$
gives $\varphi''(0)=-2$. Hence, by \eqref{eq:deltan},
\begin{equation}
\ip{\delta}{\varphi}=1,\qquad
\ip{\delta'}{\varphi}=-1,\qquad
\ip{\delta''}{\varphi}=-2 .
\end{equation}
\end{proof}

\begin{exercise}\label{ex4}
Compute $\F[3-2t+t^{2}]$ and verify the symmetry statement of
Proposition~\ref{prop:symmetries} in distributional form.
\end{exercise}

\begin{proof}[Solution]
By Theorem~\ref{thm:tn} and linearity,
\begin{equation}
\F[3-2t+t^{2}]=6\pi\delta-4\pi i\,\delta'-2\pi\,\delta'' .
\end{equation}
The function $3+t^{2}$ is real and even; it contributes $6\pi\delta-2\pi\delta''$, which is real and
even as a distribution. The term $-2t$ is real and odd and contributes $-4\pi i\delta'$, which is
purely imaginary and odd. Since $\delta^{(n)}$ is even for $n$ even and odd for $n$ odd
(it reads off $\varphi^{(n)}(0)$, and reflection changes $n$-th derivatives by $(-1)^{n}$), the
elementary symmetries are reproduced exactly.
\end{proof}

\begin{exercise}\label{ex5}
Prove $x\delta'=-\delta$ and $x\,\PV\frac{1}{x}=1$, and deduce $i\omega\hf=1$ for
$\hf=\pi\delta-i\PV\frac{1}{\omega}$.
\end{exercise}

\begin{proof}[Solution]
For the first, $\ip{x\delta'}{\varphi}=\ip{\delta'}{x\varphi}=-(x\varphi)'(0)=-\varphi(0)$.
For the second, by \eqref{eq:PVint} with $x\varphi$ in place of $\varphi$,
\begin{equation}
\left\langle x\PV\tfrac{1}{x},\varphi\right\rangle
=\int_{0}^{\infty}\frac{x\varphi(x)+x\varphi(-x)}{x}\dd x
=\int_{0}^{\infty}\big(\varphi(x)+\varphi(-x)\big)\dd x=\int_{\R}\varphi .
\end{equation}
Applying multiplication by $i\omega$ to the candidate spectrum,
\begin{equation}
i\omega\Big[\pi\delta-i\PV\tfrac{1}{\omega}\Big]
=i\pi\underbrace{\omega\delta}_{=0}+\underbrace{\omega\PV\tfrac{1}{\omega}}_{=1}=1 ,
\end{equation}
which is \eqref{eq:iomegaH}. This also displays the ambiguity: adding $c\delta$ changes nothing.
\end{proof}

\begin{exercise}\label{ex6}
Show that the value assigned to $H(0)$ does not affect the distribution it defines.
\end{exercise}

\begin{proof}[Solution]
Let $H_{1},H_{2}$ differ only at $t=0$, with $|H_{1}(0)-H_{2}(0)|=c<\infty$. For
$\varphi\in C_{c}^{\infty}$,
\begin{equation}
\ip{H_{1}}{\varphi}-\ip{H_{2}}{\varphi}
=\int_{\R}(H_{1}-H_{2})\varphi\dd t=0 ,
\end{equation}
because the integrand vanishes except at a single point and the Lebesgue integral ignores null sets.
\end{proof}

\begin{exercise}\label{ex7}
Explain why no single $\Lone$ function dominates the family $K_{\ep}$, and why the substitution
$x=\ep u$ repairs the situation.
\end{exercise}

\begin{proof}[Solution]
If $K_{\ep}\le g\in\Lone$ for all $\ep$, then $K_{\ep}\to0$ pointwise off the origin and dominated
convergence gives $\int K_{\ep}\to0$, contradicting $\int K_{\ep}=1$.

The substitution does not produce a dominator for $K_{\ep}$; it changes the object being integrated:
\begin{equation}
\int_{\R}K_{\ep}(x)\varphi(x)\dd x=\frac{1}{\pi}\int_{\R}\frac{\varphi(\ep u)}{1+u^{2}}\dd u ,
\end{equation}
and now the integrand is dominated by $\|\varphi\|_{\infty}/(\pi(1+u^{2}))$, independent of $\ep$ and
integrable. The varying kernel has been replaced by a fixed weight, at the cost of letting the test
function depend on $\ep$.
\end{proof}

\begin{exercise}\label{ex8}
Exhibit a family with unit mass whose distributional limit is $0$, and reconcile this with mass
conservation.
\end{exercise}

\begin{proof}[Solution]
Take $k_{n}=\tfrac12\mathbf{1}_{[n-1,n+1]}$. For $\varphi\in\Sw$,
$|\ip{k_{n}}{\varphi}|\le\tfrac12\int_{n-1}^{n+1}|\varphi|=O(n^{-2})\to0$. There is no contradiction:
$\int k_{n}=1$ is a statement about the functional $\varphi\equiv1$, and $1\notin\Sw$. Total mass is
therefore not a continuous functional on $\Sw'$; the mass is conserved but escapes to infinity, where
no rapidly decaying probe can detect it.
\end{proof}

\begin{exercise}\label{ex9}
Using only the dual definition, prove $\F[\delta]=1$ and $\F[1]=2\pi\delta$, and identify where
inversion enters.
\end{exercise}

\begin{proof}[Solution]
\begin{equation}
\ip{\F\delta}{\varphi}=\ip{\delta}{\ph}=\ph(0)=\int_{\R}\varphi=\ip{1}{\varphi},
\end{equation}
because $\ph(0)=\int\varphi(t)\dd t$. For the second,
\begin{equation}
\ip{\F1}{\varphi}=\ip{1}{\ph}=\int_{\R}\ph(t)\dd t=2\pi\varphi(0)=\ip{2\pi\delta}{\varphi},
\end{equation}
and the middle equality is exactly \eqref{eq:inversion} evaluated at $t=0$.
\end{proof}

\begin{exercise}\label{ex10}
Compute $\F[H(\cdot-a)]$ and explain why the $\delta$ term is unaffected by the shift.
\end{exercise}

\begin{proof}[Solution]
By \eqref{eq:timeshift},
\begin{equation}
\F[H(\cdot-a)](\omega)=e^{-ia\omega}\Big[\pi\delta(\omega)-i\PV\frac{1}{\omega}\Big]
=\pi\delta(\omega)-i\,e^{-ia\omega}\PV\frac{1}{\omega}.
\end{equation}
The factor $e^{-ia\omega}$ is a smooth multiplier and $e^{-ia\cdot0}=1$, so by \eqref{eq:adelta}
\begin{equation}
e^{-ia\omega}\,\delta(\omega)=\delta(\omega).
\end{equation}
The principal value term keeps the factor, because $\PV(1/\omega)$ is not concentrated at a point.
\end{proof}

\begin{exercise}\label{ex11}
Locate the two sign changes that produce the $+i$ in $\F(tT)=i\,\partial_{\omega}\widehat T$.
\end{exercise}

\begin{proof}[Solution]
First, from $\F[\varphi']=i\omega\ph(\omega)$ with the variable renamed,
$t\ph(t)=-i(\varphi')^{\wedge}(t)$; the factor $1/i=-i$ is the first sign. Second, the definition of
the distributional derivative gives
$\ip{\widehat T'}{\varphi}=-\ip{\widehat T}{\varphi'}$, a second minus. Then
\begin{equation}
\ip{\F(tT)}{\varphi}=\ip{T}{t\ph}
=-i\ip{T}{(\varphi')^{\wedge}}
=-i\ip{\widehat T}{\varphi'}
=i\ip{\widehat T'}{\varphi}.
\end{equation}
The two negatives cancel, leaving $+i$.
\end{proof}

\begin{exercise}\label{ex12}
For $K_{\ep}'(x)=-2\ep x/(\pi(x^{2}+\ep^{2})^{2})$ determine the location and size of the extrema,
and compute $\int|K_{\ep}'|$.
\end{exercise}

\begin{proof}[Solution]
Differentiating $\log|K_{\ep}'|$ gives the condition
$\frac{\dd}{\dd x}\big(\log x-2\log(x^{2}+\ep^{2})\big)=1/x-4x/(x^{2}+\ep^{2})=0$, that is
$\ep^{2}-3x^{2}=0$, so $x=\pm\ep/\sqrt3$. Substituting,
\begin{equation}
|K_{\ep}'(\ep/\sqrt3)|
=\frac{2\ep\cdot\ep/\sqrt3}{\pi\big(\ep^{2}+\ep^{2}/3\big)^{2}}
=\frac{2\ep^{2}/\sqrt3}{\pi\cdot 16\ep^{4}/9}
=\frac{9}{8\sqrt3\,\pi\ep^{2}} .
\end{equation}
For the total variation, use symmetry and the substitution $u=x^{2}+\ep^{2}$:
\begin{equation}
\int_{\R}|K_{\ep}'|\dd x
=2\int_{0}^{\infty}\frac{2\ep x}{\pi(x^{2}+\ep^{2})^{2}}\dd x
=\frac{4\ep}{\pi}\cdot\frac{1}{2\ep^{2}}
=\frac{2}{\pi\ep}.
\end{equation}
Since $\int K_{\ep}'=0$, the positive and negative masses are each $1/(\pi\ep)$, both diverging.
\end{proof}

\begin{exercise}\label{ex13}
Show that $e^{x^{2}}$ defines a distribution but not a tempered distribution.
\end{exercise}

\begin{proof}[Solution]
Local integrability is clear, so Theorem~\ref{thm:embedding} applies. Suppose it were tempered, so
that \eqref{eq:Sprime-est} held. Fix $0\not\equiv\psi\in C_{c}^{\infty}$ nonnegative and set
$\psi_{R}(x)=\psi(x-R)$. Then $p_{m,n}(\psi_{R})=O(R^{m})$ for each fixed $m,n$, so the right side of
\eqref{eq:Sprime-est} is $O(R^{M})$, whereas
\begin{equation}
\ip{e^{x^{2}}}{\psi_{R}}\ge e^{(R-r)^{2}}\int\psi
\end{equation}
grows faster than any polynomial. Contradiction.
\end{proof}

\begin{exercise}\label{ex14}
Prove the second identity of Lemma~\ref{lem:two-identities}, namely
$\omega^{m}\ph(\omega)=i^{-m}\F[\varphi^{(m)}](\omega)$.
\end{exercise}

\begin{proof}[Solution]
By Theorem~\ref{thm:derivt} applied to $\varphi^{(m)}$,
$\F[\varphi^{(m)}]=(i\omega)^{m}\ph(\omega)$. Multiplying both sides by $i^{-m}$ and using
$i^{-m}i^{m}=1$ gives the claim.
\end{proof}

\begin{exercise}\label{ex15}
Show that if $f$ is real and even then $\fh$ is real and even, and verify this for
$G_{\ep}=\F[e^{-\ep|\cdot|}]$.
\end{exercise}

\begin{proof}[Solution]
Evenness of $f$ gives, after $t\mapsto-t$,
\begin{equation}
\fh(-\omega)=\int f(t)e^{i\omega t}\dd t=\int f(-s)e^{-i\omega s}\dd s=\int f(s)e^{-i\omega s}\dd s=\fh(\omega),
\end{equation}
and $f$ real gives $\overline{\fh(\omega)}=\int f(t)e^{i\omega t}\dd t=\fh(-\omega)$, hence
$\fh=\overline{\fh}$. For $f_{\ep}=e^{-\ep|t|}$, which is real and even, one gets
$G_{\ep}=2\ep/(\ep^{2}+\omega^{2})$, visibly real and even.
\end{proof}

\begin{exercise}\label{ex16}
Where exactly does the proof of Theorem~\ref{thm:inversion} use Fubini, and what makes the double
integral absolutely convergent?
\end{exercise}

\begin{proof}[Solution]
Fubini is used to pass from
$\frac{1}{2\pi}\iint\varphi(s)e^{-i\omega s}e^{i\omega t}e^{-\eta\omega^{2}/2}\dd s\dd\omega$ to the
iterated integral with the inner integral over $\omega$. Absolute convergence holds because the modulus
of the integrand factors as $|\varphi(s)|\,e^{-\eta\omega^{2}/2}$, and the product of two integrable
functions is integrable on the product space: $\varphi\in\Lone$ since $\varphi\in\Sw$, and
$e^{-\eta\omega^{2}/2}$ is integrable in $\omega$. The Gaussian factor is indispensable; without it the
double integral is not absolutely convergent and the interchange is unjustified.
\end{proof}


\end{document}
